% concatenated from ldg_conv_min.tex
\documentclass[a4paper]{amsart}

%\usepackage[T1]{fontenc}
\usepackage[utf8]{inputenc}
\usepackage[british]{babel}
\usepackage{amsfonts,amsmath,amssymb,amsthm,mathtools,esint}
\usepackage{tikz}
\usepackage{enumitem}
\usepackage[hidelinks]{hyperref}
\usepackage{cleveref}
\crefname{subsection}{subsection}{subsections}
\usepackage{thmtools}
\numberwithin{equation}{section}
\usepackage{xcolor}
\definecolor{dark-gray}{gray}{0.3}
\usepackage{caption}
\setlength{\textfloatsep}{9pt}

%geometry
\newcommand{\dist}{\mathrm{dist}}
\newcommand{\diam}{\mathrm{diam}}
\newcommand{\R}{\mathbb{R}}
\newcommand{\Sym}{\mathbb{S}}
\renewcommand{\d}[1]{\,\mathrm{d}#1}
\renewcommand{\div}{\mathrm{div}}
\newcommand{\I}{\mathrm{I}}
\newcommand{\vertiii}[1]
{{\left\vert\kern-0.25ex\left\vert\kern-0.25ex\left\vert #1 
\right\vert\kern-0.25ex\right\vert\kern-0.25ex\right\vert}}

% triangulation
\newcommand{\Fcal}{\mathcal{F}}
\newcommand{\Tcal}{\mathcal{T}}
\newcommand{\tr}{\mathrm{tr}}
\newcommand{\Mcal}{\mathcal{M}}
\newcommand{\Acal}{\mathcal{A}}

% HHO
\newcommand{\PotRec}{\mathcal{R}}
\newcommand{\RT}{\mathrm{RT}}
\newcommand{\pw}{\mathrm{pw}}
\newcommand{\D}{\mathrm{D}}
\newcommand{\s}{\mathrm{s}}
\newcommand{\M}{\mathbb{M}}
\newcommand{\osc}{\mathrm{osc}}

% theorems
\newtheorem{theorem}{Theorem}[section]
\newtheorem{lemma}[theorem]{Lemma}
\newtheorem{definition}[theorem]{Definition}
\newtheorem{corollary}[theorem]{Corollary}
\newtheorem{proposition}[theorem]{Proposition}
\theoremstyle{remark}
\newtheorem{remark}[theorem]{Remark}
\newtheorem{example}[theorem]{Example}

%%% Local Variables:
\newcounter{cntS}
\newcommand{\newcnstS}{%
	\refstepcounter{cntS}%
	\ensuremath{c_{\thecntS}}}
\newcommand{\cnstS}[1]{\ensuremath{c_{\ref{#1}}}}

\newcounter{cntL}
\newcommand{\newcnstL}{%
	\refstepcounter{cntL}%
	\ensuremath{C_{\thecntL}}}
\newcommand{\cnstL}[1]{\ensuremath{C_{\ref{#1}}}}

\usepackage[backend=bibtex, giveninits=true, isbn=false, doi=false, url=false, maxnames = 4]{biblatex}
\DeclareFieldFormat[article]{title}{{#1}} % Remove "..." for articles in references
\DeclareFieldFormat[book]{title}{{#1}}
\DeclareFieldFormat[inproceedings]{title}{{#1}}
\DeclareFieldFormat[incollection]{title}{{#1}}
\DeclareFieldFormat{pages}{{#1}} % Remove pp. from pp. 101-123 in references
\AtBeginBibliography{\footnotesize}
\addbibresource{references.bib}

\begin{document}
\title[LDG for convex minimization]{Local discontinuous Galerkin \\ FEM for convex minimization}
\author[C.~Carstensen, N.~T.~Tran]{Carsten Carstensen \& Ngoc Tien Tran}
\thanks{The second author received funding from the European Union's Horizon 2020 
research and innovation programme (project RandomMultiScales, grant agreement No.~865751).}
\address[C.~Carstensen]{%     
	Humboldt-Uni\-ver\-si\-t\"at zu Berlin,
	10117 Berlin, Germany}
\email{cc@math.hu-berlin.de}
\address[N.~T.~Tran]{Universit\"at Augsburg, 86159 Augsburg, Germany}
\email{ngoc1.tran@uni-a.de}
%\date{\today}
\keywords{discrete convex duality, local discontinuous Galerkin, hybridizable method, convex minimization, a~priori, a~posteriori}
\subjclass{65N12, 65N30, 65Y20}
\begin{abstract}
The heart of the a priori and a posteriori error control in convex minimization problems 
    is the sharp control of the differences of discrete and exact minimal 
    energy. Conforming finite element discretizations  for p-Laplace type minimization problems
    provide upper bounds of the energy difference with optimal convergence rates.
    Even for smooth solutions, known convergence rates for higher-order non-conforming 
    finite element discretizations for the same problem class with $2 < p < \infty$,  however,  are exclusively 
    suboptimal. Thus the popular a posteriori 
    error control within the two-energy principle, that generalize hyper-circle identities, 
    appears unbalanced. 
     
The innovative point of departure in a refined analysis of two discontinuous Galerkin 
    (dG) schemes exploits duality relations between a discrete 
    primal and a semi-discrete dual problem. The infinite-dimensional dual problem
    leads to a tiny duality gap that even vanishes for polynomial low-order terms.
    For a class of degenerated convex minimization problems with two-sided $p$ growth,
    the novel duality
    provides improved a priori convergence rates for the error in the minimal energies. 
    % Einschub CC
    This closes the misfit of convergence rates for the conforming and nonconforming 
    schemes at least for the local discontinuous Galerkin schemes at hand.
    %
    The  motivating two-energy principle and some post-processing for a Raviart-Thomas
    dual variable provides an a posteriori error control, that also 
    may drive adaptive mesh-refining. Computational benchmarks  provide striking   
    numerical evidence for improved convergence rates of the adaptive beyond uniform
    mesh-refining.
\end{abstract}

\maketitle

\section{Introduction}\label{sec:main-results}
This paper develops novel techniques to establish duality relations for higher-order nonconforming methods for convex minimization problems with application to the error analysis of local discontinuous Galerkin methods (LDG).

\subsection{Model problem}\label{sec:continuous-problem}
Given an open bounded polyhedral Lipschitz domain $\Omega \subset \mathbb{R}^n$, the continuous problem minimizes the energy
\begin{align}
	E(v) \coloneqq \int_\Omega (W(\nabla v) + \psi(\bullet,v)) \d{x} \quad\text{among } v \in V \coloneqq W^{1,p}_\mathrm{D}(\Omega)\label{def:energy}
\end{align}
with  a convex energy density $W \in C(\mathbb{R}^n)$ and with 
a measurable function $\psi : \Omega \times \mathbb{R} \to \mathbb{R} \cup \{+\infty\}$,  convex
in the second variable: At a.e. $x \in \Omega$ let 
$\psi(x, \bullet):\R\to\R $ be a  proper lower semi-continuous convex function. 
Underlying non-displayed growth conditions lead to 
$1<p<\infty$  and the test space 
$V = W^{1,p}_\mathrm{D}(\Omega)$ of Sobolev functions in $W^{1,p}(\Omega)$ with homogeneous boundary data on a compact part $\Gamma_\mathrm{D} \subset \partial \Omega$ of the boundary $\partial \Omega$ with 
positive surface measure. 

The convexity in the functions inside \eqref{def:energy} gives rise to a dual energy: 
Let $W^*$ (resp.~$\psi^*(x,\bullet)$) denote the convex conjugate of $W$ (resp.~$\psi(x,\bullet)$ for a.e.~$x \in \Omega$) and 
the H\"older conjugate 
$1< p' \coloneqq p/(p-1) <\infty$ of $p$, $1/p + 1/p' = 1$.
%
The dual problem of \eqref{def:energy} maximizes the dual energy
\begin{align}\label{def:dual-energy}
	E^*(\tau) \coloneqq -\int_\Omega (W^*(\tau) + \psi^*(\bullet, \div \tau)) \d{x}
    \quad\text{among } \tau \in \Sigma \coloneqq W^{p'}_\mathrm{N}(\div,\Omega). 
\end{align}
The space $\Sigma$ consists 
of vector fields 
$\tau \in L^{p'}(\Omega)^n$ with a distributional divergence $\div\, \tau$ in  $L^{p'}(\Omega)$ 
and normal traces $\tau \cdot \nu=0$, that vanish along the 
(relativly open and possibly empty) Neumann boundary $\Gamma_\mathrm{N} \coloneqq \partial \Omega \setminus \Gamma_\mathrm{D}$. It is a Banach space
%$ $W^{p'}_\mathrm{N}(\div,\Omega)$
under the graph norm as, e.g., in the Hilbert case  $W^2_\mathrm{N}(\div,\Omega) = H_\mathrm{N}(\div,\Omega)$.
Throughout this paper, the subsequent general conditions 
\begin{enumerate}[wide]
    \item[(A1)] $-\infty < \inf E(V) \leq E(v_0) < \infty$ for some $v_0 \in V$ and $E$ is continuous at $v_0$,
    \item[(A2)] $E(v) \to \infty$ as $\|\nabla v\|_p \to \infty$ for $v \in V$
\end{enumerate}
first  imply existence of solutions  $u \in \arg\min E(V)$ 
and $\sigma \in \arg \max E^*(\Sigma)$), as well as, second,  no duality gap 
\cite[Chapter 3, Theorem 4.2]{EkelandTeman1999} viz.
\begin{align}\label{eq:strong-duality}
    E(u) = \min E(V) = \max E^*(\Sigma) = E^*(\sigma).
\end{align}
 
\subsection{Motivation}
The a~priori error analysis of conforming finite element methods (FEM)
is well understood in the literature \cite{GlowinskiMarrocco1975,Chow1989,CPlechac1997}. 
We illustrate the main results on the a priori convergence rates from a class
of problems with $p \geq 2$ and generic constants $c_1,c_3,c_5 > 0$ and $c_2,c_4 \geq 0$ in the subsequent assumptions. 
\begin{enumerate}[wide]
    \item[(B1)] (smoothness of $W$) $W \in C^1(\mathbb{R}^n)$.
    \item[(B2)] (two-sided growth) Any $a \in \mathbb{R}^n$ satisfies
$       c_1|a|^p - c_2 \leq W(a) \leq c_3|a|^p + c_4  $.
    \item[(B3)] (convexity control)
Any $a,b \in \mathbb{R}^n$ satisfy\begin{align*}
    d(a,b) \coloneqq \frac{ |\D W(a) - \D W(b)|^2}{ 1 + |a|^{p-2} + |b|^{p-2}}
\leq c_5\left(W(b) - W(a) - \D W(a) \cdot (b - a)\right).
\end{align*}   % $d(a,b) \coloneqq |\D W(a) - \D W(b)|^2/(1 + |a|^{p-2} + |b|^{p-2})$.
\item[(B4)] (linear low-order term) $\psi(x,a) \coloneqq -f(x) a$ 
for some $f \in L^{p'}(\Omega)$.
\end{enumerate}
The assumptions (B1)--(B3) define a class of degenerate convex minimization problems:
The dual variable $\sigma = \D W(\nabla u)$ is uniquely defined, i.e.,  
independent of the choice of the possibly non-unique minimizer $u$, 
and $\sigma$ belongs to  $H^1_{\rm loc}(\Omega)^n$ 
\cite{CPlechac1997,CarstensenMueller2002}.
%
Examples include an optimal design problem \cite{KohnStrang1986,BartelsC2008} and relaxed 
scalar double well problems \cite{CPlechac1997}.
%
The condition (B4) allows for explicit error control, but the analysis of this paper extends to other right-hand sides as well.

Given a discrete minimizer $u_h \in \arg\min E(V_h)$ of the energy 
\eqref{def:energy}
in a conforming subspace  $V_h \subset V$,  let $\sigma_h \coloneqq \D W(\nabla u_h)$ denote the discrete dual variable. 
Under the assumptions (B1)--(B4),
arguments from \cite{GlowinskiMarrocco1975,Chow1989,CPlechac1997} imply 
\begin{align}\label{confenergyerror}
E(u_h) - E(u) \lesssim \inf_{v_h \in V_h} \|\nabla(u - v_h)\|_p^2 
=O(h_{\max}^{2k})
\end{align}
with the maximal mesh-size $h_{\max}$ of the underlying finite element mesh
and under sufficient smoothness assumptions 
with the discretization order $k \geq 1$. For the convenience of the reader, we provide a proof of \eqref{confenergyerror} in the appendix.
This control of the error in the energies \eqref{confenergyerror}  enables several
error estimates for the error in the primal and dual variables
in \cite{GlowinskiMarrocco1975,Chow1989,CPlechac1997}.

The a priori error analysis of higher order nonconforming schemes is less developed
and we cannot even quote directly an analog of \eqref{confenergyerror} from the literature. Before
we present  \eqref{nonconfenergyerror},
we point out  that known  a priori error estimates for higher-order nonconforming schemes are  
really suboptimal:  For strongly monotone problems and $p\ge 2$,
\cite{DiPietroDroniou2017-II,DroniouEymardGallouet2018} 
(and the references therein) solely 
establish 
$\|\nabla_\pw(u-u_h)\|_{L^p(\Omega)} =O(h_{\max}^{k/(p-1)})$,
which is inferior to \eqref{confenergyerror}
%$\|\nabla(u-u_h)\|_{L^p(\Omega)} = O(h_{\max}^{2k/p})$  
for the
conforming FEM with \eqref{confenergyerror} \cite{Chow1989}.

A similar sub-optimality  arises in the case $1 < p < 2$, which has been resolved only recently in \cite{Tran2024} with weak duality 
between HHO and discrete dual hybrid methods.
In the current case $p\ge 2$, however, those
arguments cannot fully mimic the proof of 
\eqref{confenergyerror} and led  to suboptimal rates.

The main motivation to apply nonconforming methods comes from a direct 
a~posteriori error control by duality 
\cite{Repin1997,Repin2000,NeittaanmakiRepin2004,CLiu2015,Bartels2015} in 
degenerate problems, 
when residual-based error estimation suffer from the reliability-efficiency gap \cite{CJochimsen2003}.


The smoothness assumptions for  \eqref{confenergyerror}-\eqref{nonconfenergyerror} may appear unrealistically high,
but behind those short statements is an error analysis that reduces the convergence of the schemes to that of
interpolation errors. On the theoretical side, optimal rates indicate a sharp analysis. On the practical side
we might expect that an adaptive mesh-refining miraculously  removes singularity-driven suboptimal global 
approximation errors: Computational benchmarks may confirm this wishful vision.


\subsection{Contributions of this paper}
% It is well understood in the literature that certain classes of discontinuous Galerkin methods can be hybridized for linear problems \cite{CockburnGopalakrishnanLazarov2009,DiPietroErn2012} by setting appropriate numerical traces. This concept was also observed for problems of $p$-Laplace type \cite{LDG2014}. 
This paper establishes the analog of \eqref{confenergyerror} for 
local discontinuous Galerkin (LDG) methods %in the spirit of 
\cite{CockburnShu1998,BurmanErn2008,DiPietroErn2012}, viz.
\begin{align}\label{nonconfenergyerror}
E_h(\mathrm{I}_h u) - \min E_h(V_h) + |E(u) - \min E_h(V_h)|
=O(h_{\max}^{2k}),
\end{align}
with novel strong duality relations on the discrete level by employing continuous objects in the dual ansatz space.
%to overcome the non-linearity of the energy density for higher-order discretizations.
% Furthermore, it turns out that this equivalence holds between HHO/HDG methods with the Lehrenfeld-Sch\"oberl \cite{lehrenfeld2010hybrid} as well as the early suboptimal WG stabilization. Therefore, these hybridizable and LDG methods are equivalent and this allows for a characterization of the trace variables of the hybridizable methods considered in this paper.
This leads in the a~priori error analysis to %the estimate 
\begin{align*}
	0 \leq E_h(\mathrm{I}_h u) - \min E_h(V_h) \leq E_h(\mathrm{I}_h u) - E(u) + \gamma_h(\mathrm{I}_h^* \sigma)/r'.
\end{align*}
Here and in \eqref{nonconfenergyerror}, 
$E_h$ denotes a dicrete version of \eqref{def:energy}, 
$V_h$ is the discrete ansatz space, 
and $\mathrm{I}_h$ and $\mathrm{I}_h^*$ are interpolation operators. 
Up to the stabilization $\gamma_h(\mathrm{I}_h^* \sigma)/r'$ on the dual level, the discrete energy of the interpolation of $u$ is an {\em upper bound} for the exact energy. Under the assumptions (B1)-(B4) and
% (possibly unrealistic) 
piecewise 
smoothness assumptions $u \in W^{k+1,p}$ and $\sigma \in W^{k,2}$, we provide quadratic convergence rates \eqref{nonconfenergyerror}.
This enables 
error estimates in the primal and dual variables as in \cite[Section 4]{Tran2024}, e.g.,
$\|\nabla_\pw(u-u_h)\|_{p} = O(h_{\max}^{2k/p})$ follows
for $p \geq 2$ in strongly monotone problems. 
%
This closes the theoretical gap between conforming and nonconforming discretizations.

%WAS IST DER STATUS HIER GENAUER?
Our results carry over to hybridizable methods with  HDG/WG stabilizations \cite{CockburnGopalakrishnanLazarov2009} with suboptimal polynomial consistency ~\cite[Remark 2.9]{DiPietroDroniou2017}; however, it does not apply 
for Lehrenfeld-Sch\"oberl stabilization.

A conforming dual  Raviart-Thomas finite element functions  approximation 
enables guaranteed energy error bounds. We suggest a post-processing by
explicit design of the required degrees of freedom as in \cite{ErnStephansenVohralik2010} for linear problems. 
A localization of the resulting error estimator drives an 
adaptive mesh-refining algorithms as an alternative to \cite{CarstensenTran2021}. 

\subsection{Outline}
The remaining parts of this paper are organized as follows. \Cref{sec:discretization} introduces the numerical methods considered in this paper. The equivalence of these methods to dual maximization problems is established in \Cref{sec:equivalence}. This applies to the error analysis of an LDG method in \Cref{sec:error-analysis}. 
An extension of the analysis to hybridizable methods is briefly discussed in \Cref{sec:hybrid}.
Three numerical benchmarks in \Cref{sec:numerical_examples} 
with improved convergence rates for adaptive mesh-refining algorithms
conclude this paper.  

\subsection{Notation}
Standard notation for Sobolev and Lebesgue spaces applies throughout this paper with the abbreviation $\|\bullet\|_p \coloneqq \|\bullet\|_{L^p(\Omega)}$ for any $1 < p < \infty$.
The notation $A \lesssim B$ abbreviates $A \leq CB$ for a generic constant $C$ independent of the mesh-size and $A \approx B$ abbreviates $A \lesssim B \lesssim A$.

\section{Discretization}\label{sec:discretization}
This section presents the numerical scheme for the discretization of \eqref{def:energy}.

\subsection{Polytopal Mesh}\label{sec:triangulation}
Let $\mathcal{M}$ be a finite collection of closed polytopes of positive volume with overlap of measure zero that covers $\overline{\Omega} = \cup_{K \in \mathcal{M}} K$.
A face $S$ of the mesh $\mathcal{M}$ is a closed connected subset of a hyperplane $H_S$ with positive $(n-1)$-dimensional surface measure such that either (a) there exist $K_+,K_- \in \mathcal{M}$ with $S \subset H_S \cap K_+ \cap K_-$ (interior face) or (b) there exists $K_+ \in \mathcal{M}$ with $S \subset H_S \cap K_+ \cap \partial \Omega$ (boundary face). We refer to \cite[Section 1.1]{DiPietroDroniou2017} for further details.

Let $\mathcal{F}$ be a finite collection of faces with overlap of $(n-1)$-dimensional surface measure zero that covers the skeleton $\partial \Mcal \coloneqq \cup_{K \in \mathcal{M}} \partial K = \cup_{S \in \mathcal{F}} S$ with the split $\mathcal{F} = \mathcal{F}(\Omega) \cup \mathcal{F}(\partial \Omega)$ into the set of interior faces $\mathcal{F}(\Omega)$ and the set of boundary faces $\mathcal{F}(\partial \Omega)$.
Let $\mathcal{F}_\mathrm{D} \coloneqq \{S \in \mathcal{F} : S \subset \Gamma_\mathrm{D}\}$ (resp.~$\mathcal{F}_\mathrm{N} \coloneqq \mathcal{F}(\partial\Omega) \setminus \mathcal{F}_\mathrm{D}$) denote the set of Dirichlet (resp.~Neumann) faces. For $K \in \mathcal{M}$, $\mathcal{F}(K)$ is the set of all faces of $K$.
The normal vector $\nu_S$ of an interior face $S \in \mathcal{F}(\Omega)$ is fixed in its orientation beforehand and set $\nu_S \coloneqq \nu|_S$ for boundary faces $S \in \mathcal{F}(\partial \Omega)$.
For $S \in \mathcal{F}(\Omega)$, $K_+ \in \Mcal_h$ (resp.~$K_- \in \Mcal$) denotes the unique cell with $S \subset \partial K_{+}$ (resp.~$S \subset \partial K_-$) and $\nu_{K_+}|_S = \nu_S$ (resp.~$\nu_{K_-}|_S = -\nu_S$).
For $S \in \mathcal{F}(\partial \Omega)$, $K_+ \in \Mcal$ is the unique cell with $S \subset \partial K_+$.
The jump $[v]_S$ and the average $\{v\}_S$ of any function $v \in W^{1,1}(\mathrm{int}(T_+ \cup T_-))$ along $S \in \mathcal{F}(\Omega)$ are defined by $[v]_S \coloneqq v|_{K_+} - v|_{K_-} \in L^1(S)$ and $\{v\}_S \coloneqq (v|_{K_+} + v|_{K_-})/2 \in L^1(S)$.
If $S \in \mathcal{F}(\partial \Omega)$, then $[v]_S \coloneqq v|_S = v_{K_+}|_S \eqqcolon \{v\}_S$.

For theoretical purposes, let $\vartheta$ denote the mesh regularity parameter of $\Mcal$ associated with a matching simplicial submesh, we refer to \cite[Definition 1.38]{DiPietroErn2012} for a detailed definition. The constants in discrete inequalities such as the trace or inverse inequality depend on this parameter.
%The notation $c_\mathrm{P} \coloneqq \max_{K \in \Mcal} c_\mathrm{P}(K)$ denotes the best possible local Poincar\'e constant on $\Mcal$. For general polyhedral meshes, $c_\mathrm{P}$ may depend on $\vartheta$.
The differential operators $\nabla_\pw$ and $\div_\pw$ denote the piecewise version of $\nabla$ and $\div$ without explicit reference to the underlying mesh.

\subsection{Finite element spaces}\label{sec:fem-spaces}
Given a subset $M \subset \R^n$ of diameter $h_M$, let $P_k(M)$ denote the space of polynomials of degree at most $k$.
For any $v \in L^1(M)$, $\Pi_M^k v \in P_k(M)$ denotes the $L^2$ projection of $v$ onto $P_k(M)$;
$P_k(\mathcal{M})$ and $P_k(\mathcal{F})$ denote 
the space of piecewise polynomials of degree at most $k$ with respect to the mesh $\mathcal{M}$ and the faces $\mathcal{F}$.
%Given $v \in L^1(\Omega)$, we define the $L^2$ projection $\Pi_\Mcal^k v$ of $v$ onto $P_k(\Mcal)$ by $(\Pi_\Mcal^k v)|_K = \Pi_K^k v|_K$ in any cell $K \in \Mcal$.
%We recall the stability of the $L^2$ projection $\|\Pi^k_\Mcal \xi\|_p \lesssim \|\xi\|_p$ for any $\xi \in L^p(\Omega)$
%from, e.g., \cite[Lemma 3.2]{DiPietroDroniou2017}. 
%A consequence of this is the quasi-best approximation
%\begin{align}\label{ineq:quasi-opt-L2}
%	\|(1 - \Pi^k_\Mcal) \xi\|_p \leq \|\xi - \varphi_k\|_p + \|\Pi^k_\Mcal(\varphi_k - \xi)\|_p \lesssim \|\xi - \varphi_k\|_p
%\end{align}
%for any $\varphi_{k} \in P_k(\Mcal)$, where $\varphi_{k} = \Pi_\Mcal^k \varphi_k$ is utilized in the penultimate inequality.
The piecewise constant function $h_\Mcal \in P_0(\Mcal)$ reads $h_\Mcal|_K = h_K = \mathrm{diam}(K)$; $h_{\max} \coloneqq \max_{K \in \Mcal} h_K$ is the maximal mesh-size of $\Mcal$.
%For any $f \in L^{p'}(\Omega)^m$, the oscillation of $f$ reads $\osc(f) \coloneqq \|h_\Mcal(1 - \Pi_\Mcal^{k+1}) f\|_{p'}$.

% \subsection{Hybridizable method}\label{sec:hho}
% Given a fixed $k \in \mathbb{N}$, let $V_h \coloneqq P_{k+1}(\Mcal) \times P_\ell(\mathcal{F}\setminus\mathcal{F}_\mathrm{D})$ with $\ell \in \{k,k+1\}$ denote the discrete ansatz space. The discrete gradient $\nabla_h v_h \in P_k(\Mcal)^n$ of $v_h = (v_\Mcal, v_\mathcal{F}) \in V_h$ is the unique solution to
% \begin{align}\label{def:discrete-gradient}
% 	\int_\Omega \nabla_h v_h \cdot \Phi \d{x} = - \int_\Omega v_\Mcal \div_\pw \Phi \d{x} + \sum_{S \in \mathcal{F}\setminus\mathcal{F}_\mathrm{D}} \int_S v_S [\Phi]_S \cdot \nu_S \d{s}
% \end{align}
% for any $\Phi \in P_k(\Mcal)^n$, where the convention $v_S \coloneqq v_\mathcal{F}|_S$ (and late $v_K \coloneqq v_\Mcal|_K$) is employed throughout this paper. Let $\psi_h : \Omega \times \mathbb{R}^m \to \mathbb{R} \cup \{+\infty\}$ be a suitable approximation of $\psi$ (so that $\psi_h(x,\bullet)$ is a proper lower semicontinuous convex function for a.e.~$x \in \Omega$). Given fixed parameters $r \in (1,\infty)$ and $s \in \mathbb{R}$,
% the discrete hybridizable problem minimizes
% \begin{align}\label{def:hho_discrete_energy}
% 	E_h(v_h) \coloneqq \int_\Omega (W(\nabla_h v_h) + \psi_h(x,v_\Mcal)) \d{x} + s_h(v_h)/r
% \end{align}
% among $v_h = (v_\Mcal, v_\mathcal{F}) \in V_h$ with the stabilization
% \begin{align*}
% 	s_h(v_h;w_h) \coloneqq \sum_{K \in \Mcal} \sum_{S \in \mathcal{F}(K)} h_S^{-s} \int_S |v_S - \Pi_S^\ell v_K|^{r-2}(v_S - \Pi_S^\ell v_K)(w_S - \Pi_S^\ell w_K) \d{s}
% \end{align*}
% for any $w_h = (w_\Mcal, w_\mathcal{F}) \in V_h$. If $\ell = k$, $s_h$ is the Lehrenfeld-Sch\"oberl stabilization, while the choice $\ell = k+1$ leads to a suboptimal polynomial consistency of the stabilization. 

\subsection{Modified local discontinuous Galerkin method}\label{sec:discrete-spaces}
For any $k \geq 1$,
% The main objective of this paper is to investigate the equivalence of \eqref{def:hho_discrete_energy} and the LDG method in \eqref{def:discrete_energy} below. To this end, 
we consider the discrete ansatz space 
$$V_h \coloneqq P_{k}(\mathcal{M}).$$
% In other words, $V_h^{\mathrm{dg}}$ consists of the typical cell degrees of freedom, but also trace degrees of freedom on the Neumann boundary.
% Naturally, $V_h^{\mathrm{dg}}$ can be embedded into $V_h$ by the following injective map $v_h^{\mathrm{dg}} = (v_\Mcal, v_{\mathcal{F}_\mathrm{N}}) \mapsto \mathcal{E}_h v_h^{\mathrm{dg}} \coloneqq (v_\Mcal, v_\mathcal{F}) \in V_h$ with $v_\mathcal{F}|_S = \Pi_S^\ell \{v_\Mcal\}_S$ for $S \in \mathcal{F}(\Omega)$ and $v_\mathcal{F}|_S = v_{\mathcal{F}_\mathrm{N}}|_S$ for $S \in \mathcal{F}_\mathrm{N}$.
The discrete gradient $\nabla_h v_h \in P_{k-1}(\Mcal)^n$ of $v_h \in V_h$ is the unique solution to
\begin{align}\label{def:discrete-gradient}
 	\int_\Omega \nabla_h v_h \cdot \Phi \d{x} = - \int_\Omega v_h\, \div_\pw \Phi \d{x} + \sum_{S \in \mathcal{F}\setminus\mathcal{F}_\mathrm{D}} \int_S \{v_h\}_S [\Phi]_S \cdot \nu_S \d{s}
\end{align}
for any $\Phi \in P_{k-1}(\Mcal)^n$.
An integration by parts proves
\begin{align}\label{def:ldg-dicrete-gradient}
	\int_\Omega \nabla_h v_h \cdot \Phi \d{x}
	= \int_\Omega \nabla_\pw v_h \cdot \Phi \d{x} - \sum_{S \in \mathcal{F}\setminus\mathcal{F}_\mathrm{N}} \int_S [v_h]_S \{\Phi\}_S\cdot \nu_S \d{s}.
\end{align}
Note that $\nabla_h$ coincides with the discrete gradient of \cite[Section 4.3.2]{DiPietroErn2012} and can be expressed in terms of $\nabla_\pw$ and lifting operators \cite{BrezziManziniMariniPietraRusso2000}. The error between $\nabla_\pw v_h$ and $\nabla_h v_h$ is controlled by the interior penalty stabilization.
\begin{lemma}[discrete consistency error]\label{lem:discrete-consistency-error}
	Any $v_h \in V_h$ satisfies
	\begin{align*}
		\|\nabla_\pw v_h - \nabla_h v_h\|_{p}^p \lesssim \sum_{S \in \mathcal{F} \setminus \mathcal{F}_\mathrm{N}} h_S^{1-p}\|[v_h]_S\|_{L^p(S)}^p.
	\end{align*}
\end{lemma}
\begin{proof}
	Given any $v_h \in V_h$ and $\Phi_h \in P_{k-1}(\Mcal)^n$, a discrete trace inequality in \eqref{def:ldg-dicrete-gradient} proves
	\begin{align*}
		\int_\Omega (\nabla_\pw v_h - \nabla_h v_h) \cdot \Phi_h \d{x} \lesssim \Big(\sum_{S \in \mathcal{F} \setminus \mathcal{F}_\mathrm{N}} h_S^{1-p}\|[v_h]_S\|_{L^p(S)}^p\Big)^{1/p}\|\Phi_h\|_{p'}.
	\end{align*}
	Since $\|\Pi_\Mcal^k \Phi\|_{p'} \lesssim \|\Phi\|_{p'}$ for any $\Phi \in L^{p'}(\Omega)^n$ from stability of the $L^2$ projection in $L^{p'}(\Omega)^n$ \cite[Lemma 3.2]{DiPietroDroniou2017}, the proof concludes with
	\begin{align*}
		\|\nabla_\pw v_h - \nabla_h v_h\|_{p} &= \sup_{\Phi \in L^{p'}(\Omega)^n \setminus \{0\}} \int_\Omega (\nabla_\pw v_h - \nabla_h v_h) \cdot \Pi_\Mcal^{k-1} \Phi \d{x}/\|\Phi\|_{p'}.\qedhere
	\end{align*}
\end{proof}
Let $\psi_h : \Omega \times \mathbb{R}^m \to \mathbb{R} \cup \{+\infty\}$ be an approximation of $\psi$ (so that $\psi_h(x,\bullet)$ is a proper lower semicontinuous convex function for a.e.~$x \in \Omega$).
Given fixed parameters $1<r<\infty$ and $s \in \mathbb{R}$,
the LDG method of this paper minimizes the discrete energy
\begin{align}\label{def:discrete_energy}
	E_h(v_h) &\coloneqq \int_\Omega (W(\nabla_h v_h) 
    + \psi_h(\bullet,v_h)) \d{x} + s_h(v_h)/r,\\
    s_h(v_h; w_h) &\coloneqq 
    \sum_{S \in \mathcal{F} \setminus \mathcal{F}_\mathrm{N}} 
    h_S^{-s}\int_S |[v_h]_S|^{r-2} [v_h]_S\,[w_h]_S \d{s}
\label{def:stabilization-ldg}
\end{align}
for any $v_h, w_h \in V_h$
and the convention $s_h(v_h) \coloneqq s_h(v_h;v_h)$.
In the linear case with $W(a) \coloneqq |a|^2/2$ (and $r =2$, $s = 1$), 
this leads to a local discontinuous Galerkin method \cite{CockburnShu1998}, 
cf.~also \cite[Section 4.4.2]{DiPietroErn2012}. For the $p$-Laplace problem, 
this method was proposed in \cite{BurmanErn2008} with $r = p$ and $s = p-1$.
For the existence of discrete minimizers, we assume 
corresponding discrete versions of (A1)-(A2).
% Furthermore, we assume the following discrete version of (A1)--(A2).
% \begin{enumerate}
% 	\item[(A1d)]\label{assumption1h} There exists $v_h \in V_h$ with $E_h(v_h) < \infty$ and $-\infty < \inf E_h(V_h)$.
% 	\item[(A2d)]\label{assumption2h} $E_h(v) \to \infty$ for $v_h \in V_h$, $\|\nabla_\pw v_h\|_1 + \sum_{S \in \mathcal{F} \setminus \mathcal{F}_\mathrm{N}} \|[v_h]\|_{L^1(S)} \to \infty$.
% \end{enumerate}
\begin{remark}[other DG methods]
The design of DG methods for \eqref{def:energy} is delicate because the convex energy structure may be forfeited \cite{OrtnerSueli2007,GrekasKoumatosMakridakisVikelis2025}.
The DG methods in \cite{EyckLew2006,BurmanErn2008,BuffaOrtner2009} utilized a reconstruction operator for the discretization of the continuous gradient, which preserves the convexity structure of the energy on the discrete level.
For the lowest-order discretization on regular triangulations into simplices, the DG methods of \cite{Bartels2021} provide a simple approach by approximating the continuous gradient with the piecewise one, but the analysis requires properties 
of Crouzeix-Raviart and Raviart-Thomas finite element functions.
\end{remark}

%\subsection{Discrete dual problem}
%The ansatz space is inspired from the analysis of hybrid high-order methods in \cite{Tran2024} and reads $W_h \coloneqq P_k(\Mcal)^{m \times n} \times P_{k+1}(\mathcal{F}\setminus\mathcal{F}_\mathrm{N})^m$.
%Here, $P_{k+1}(\mathcal{F}\setminus\mathcal{F}_\mathrm{N}) \subset P_{k+1}(\mathcal{F})$ consists of functions vanishing along the Neumann boundary. 
%The interpolation $\I_h^*$ maps $\tau \in W^{1,1}(\Omega)^{m \times n} \cap W$ onto $\I_h^* \tau = (\Pi_\Mcal^k \tau, (\Pi_S^{k+1} \tau\nu_S)_{S \in \mathcal{F}\setminus\mathcal{F}_\mathrm{N}}) \in W_h$.

\section{Duality relations on discrete level}\label{sec:equivalence}
The main tools for the analysis of this paper are duality relations of the discrete problem \eqref{def:discrete_energy} to a dual maximization problem with the ansatz space
\begin{align*}
	%W_h &\coloneqq L^{p'}(\Omega)^n \times P_\ell(\mathcal{F}\setminus\mathcal{F}_\mathrm{N})\quad\text{if } \ell = k,\\
	Y &\coloneqq \{(\tau_\Mcal, \tau_\mathcal{F}) \in L^{p'}(\Omega)^n \times L^2(\partial \mathcal{M}) : \tau_\mathcal{F}|_S \equiv 0 \text{ for all } S \in \mathcal{F}_\mathrm{N}\}.
\end{align*}
Given any function $\tau \in W^{1,1}(\Omega)^n \cap \Sigma$,
we define the interpolation
\begin{align*}
	\mathrm{I}_h^* \tau \coloneqq (\tau, (\tau \cdot \nu_S)_{S \in \mathcal{F}}) \in Y.
 %    \begin{cases}
	% 	(\tau, (\Pi_S^\ell \tau \cdot \nu_S)_{S \in \mathcal{F}}) &\mbox{if } \ell = k,\\
	% 	(\tau, (\tau \cdot \nu_S)_{S \in \mathcal{F}}) &\mbox{if } \ell = k+1
	% \end{cases} \in W_h
\end{align*}
%Note that the index $h$ in $W_h$ does not indicate that the space $W_h$ is finite dimensional but it points towards the duality to the minimization of $E_h$ in $V_h$.
The divergence reconstruction $\div_h \tau \in P_{k}(\Mcal)$ of $\tau = (\tau_\Mcal, \tau_\mathcal{F}) \in Y$ is the unique solution to
\begin{align}\label{def:div-rec}
	\int_\Omega \div_h \tau\, \phi \d{x} = -\int_\Omega \tau_\Mcal \cdot \nabla_\pw \phi \d{x} + \sum_{S \in \mathcal{F} \setminus \mathcal{F}_\mathrm{N}} \int_S \tau_S [\phi]_S \d{s} 
\end{align}
for any $\phi \in P_{k}(\Mcal)$.
The operator $\div_h$ is consistent in the following sense.
\begin{lemma}[consistency]\label{lem:consistency}
	% Let $\ell = k+1$.
	Any $\tau \in W^{1,1}(\Omega)^n \cap \Sigma$ satisfies
	%\begin{align*}
		$\div_h \I_h^* \tau = \Pi_\Mcal^k \div\, \tau$.
	%\end{align*} 
\end{lemma}

\begin{proof}
	The right-hand side of \eqref{def:div-rec} is equal to $(\div\, \tau, \phi)_{L^2(\Omega)}$ for any $\phi \in P_{k}(\Mcal)$, which concludes the assertion.
\end{proof}
%Note that this critical consistency no longer holds for $\ell = k$.
%Therefore, the a~priori error analysis in \Cref{sec:error-analysis} of this paper cannot be extended to that case, we refer to \cite{Tran2024} for error estimates if $\ell = k$.
For any $\tau = (\tau_\Mcal, \tau_\mathcal{F}) \in Y$, consider the following dual energy
\begin{align}\label{def:dual-discrete-energy}
	E^*_h(\tau) &\coloneqq - \int_\Omega (W^*(\tau_\Mcal) + \psi_h^*(\bullet,\div_h \tau)) \d{x} - \gamma_h(\tau)/r',\\
    \gamma_h(\tau) &\coloneqq \sum_{S \in \mathcal{F}\setminus\mathcal{F}_\mathrm{N}} h^{s/(r-1)} \|\tau_S - \{\Pi_\Mcal^{k-1} \tau_\Mcal\}_S \cdot \nu_S\|_{L^{r'}(S)}^{r'}.\label{def:dual-stab}
\end{align}
The following duality relation between the LDG method \eqref{def:discrete_energy} and the dual problem of \eqref{def:dual-discrete-energy} holds. For strong duality, we assume the following condition.
\begin{enumerate}[wide]
    \item[(B5)] (polynomial low-order term)
    $\psi_h(x,\bullet) \in C^1(\mathbb{R})$ at a.e. $x\in \Omega$ and $\partial_u \psi_h(\bullet,v_h) \in P_k(\Mcal)$ for any $v_h \in V_h = P_k(\Mcal)$.
\end{enumerate}
 

\begin{theorem}[duality of LDG]\label{thm:weak-duality}
	It holds $\sup E_h^*(Y) \leq \min E_h(V_h)$; (B1) and (B5) imply 
	%\begin{align}%\label{eq:strong-duality-ldg}
		$\max E_h^*(Y) = \min E_h(V_h)$.
	%\end{align}
\end{theorem}

\begin{proof}
	Given $\tau = (\tau_\Mcal, \tau_\mathcal{F}) \in Y$ and $v_h \in V_h$,
    \eqref{def:ldg-dicrete-gradient} implies
    \begin{align*}
        &\int_\Omega \Pi_\Mcal^{k-1} \tau_\Mcal \cdot \nabla_h v_h \d{x}\\
        &\qquad= \int_\Omega \nabla_\pw v_h \cdot \Pi_\Mcal^{k-1} \tau_\Mcal \d{x} - \sum_{S \in \mathcal{F}\setminus\mathcal{F}_\mathrm{N}} \int_S [v_h]_S \{\Pi_\Mcal^{k-1} \tau_\Mcal\}_S \cdot \nu_S \d{s}.
    \end{align*}
    Since $\Pi_\Mcal^{k-1}$ can be omitted in the integral $(\nabla_\pw v_h, \Pi_\Mcal^{k-1} \tau_\Mcal)_{L^2(\Omega)}$,
    this and the definition of the divergence reconstruction in \eqref{def:div-rec} provide
    \begin{align}\label{eq:diby-hho}
        \int_\Omega \tau_\Mcal \cdot \nabla_h v_h \d{x} &= \int_\Omega \Pi_\Mcal^{k-1} \tau_\Mcal \cdot \nabla_h v_h \d{x} = -\int_\Omega \div_h \tau \,v_h \d{x}\nonumber\\
		&\qquad+ \sum_{S \in \mathcal{F} \setminus \mathcal{F}_\mathrm{N}} \int_S [v_h]_S (\tau_S - \{\Pi_\Mcal^{k-1} \tau_\Mcal\}_S \cdot \nu_S) \d{s}.
    \end{align}
	From this, $\tau_\Mcal \cdot \nabla_h v_h \leq W(\nabla_h v_h) + W^*(\tau_\Mcal)$ and $\div_h \tau_h \, v_h \leq \psi_h(\bullet,v_h) + \psi_h^*(\bullet,\div_h \tau)$ a.e.~in $\Omega$ as well as the H\"older inequality, we deduce $\sup E_h^*(Y) \leq \min E_h(V_h)$.
    To establish equality,
	let $u_h \in \arg\min E_h(V_h)$ be a minimizer of $E_h$ in $V_h$. We define the dual variable $y = (\sigma_\Mcal, \sigma_\mathcal{F}) \in Y$ with
	\begin{align}\label{def:dual-variable}
			\sigma_\Mcal \coloneqq \D W(\nabla_h u_h) \quad\text{and}\quad
			\sigma_\mathcal{F}|_S \coloneqq \{\Pi_{\Mcal}^{k-1} \sigma_\Mcal\}_S \cdot \nu_S - h_S^{-s}|[u_h]_S|^{r-2} [u_h]_S
	\end{align}
	for any $S \in \mathcal{F}\setminus\mathcal{F}_\mathrm{N}$.
	The Euler-Lagrange equations read
	\begin{align}\label{eq:dELE}
		0 = \int_\Omega (\sigma_\Mcal \cdot \nabla_h v_h \d{x} + \partial_u \psi_h(\bullet, u_h) v_h) \d{x} + s_h(u_h;v_h)
	\end{align}
	for any $v_h \in V_h$.
	This and
	the discrete
	integration by parts formula \eqref{eq:diby-hho} imply
	\begin{align}\label{eq:proof-mldg-equivalence}
		0 = &\int_\Omega (\partial_u \psi_h(\bullet, u_h) - \div_h y) v_h \d{x}\nonumber\\
		&\qquad + \sum_{S \in \mathcal{F} \setminus \mathcal{F}_\mathrm{N}} \int_S [v_h]_S \cdot (\sigma_S - \{\Pi_\Mcal^{k-1} \sigma_\Mcal\}_S \cdot \nu_S) \d{s} + s_h(u_h;v_h).
	\end{align}
	An explicit calculation with
	the definitions of $\sigma_\mathcal{F}$ from \eqref{def:dual-variable} and the stabilization $s_h$ from \eqref{def:stabilization-ldg} proves that the final two terms on the right-hand side cancel. Therefore, $0 = \int_\Omega (\partial_u \psi_h(\bullet, u_h) - \div_h y) v_h \d{x}$ holds for any $v_h \in V_h$ from \eqref{eq:proof-mldg-equivalence}. This yields
	\begin{align}\label{eq:div_h_sigma_h}
		\div_h y = \Pi_\Mcal^k \partial_u \psi_h(\bullet,u_h) = \partial_u \psi_h(\bullet,u_h)
	\end{align}
	under the assumption (B5).
	Since $-sr' + s/(r-1) = -r$, we obtain $\gamma_h(y) = s_h(u_h)$.
	This, the identities $\sigma_\Mcal \cdot \nabla_h u_h = W(\nabla_h u_h) + W^*(\sigma_\Mcal)$ and $\div_h y\,u_h = \psi_h(\bullet,u_h) + \psi_h^*(\bullet,\div_h y)$ from $\nabla_h u_h \in \partial W^*(\sigma_\Mcal)$ and from $u_h \in \partial_u \psi_h^*(\bullet,\div_h y)$ a.e.~in $\Omega$, and \eqref{eq:dELE} with the choice $v_h = u_h$ show
	\begin{align*}
		0 = \int_\Omega (W(\nabla_h u_h) + W^*(\sigma_\Mcal) + \psi_h(x,u_h) + \psi_h^*(x,\div_h y)) \d{x} + \frac{s_h(u_h)}{r} + \frac{\gamma_h(y)}{r'}&.
	\end{align*}
	Rearranging the terms on the right-hand side concludes the proof. 
\end{proof}

\section{Error analysis of LDG method}\label{sec:error-analysis}
In this section, we apply the duality relations in \Cref{sec:equivalence} to the error analysis. To establish error estimates, we assume for simplicity the explicit representation (B4) of the lower-order term. In this case,
\begin{align}\label{def:lower-order-terms}
	\psi_h(x,a) \coloneqq -f_h(x)\, a
\end{align}
with the $L^2$ orthogonal projection $f_h \coloneqq \Pi_\Mcal^{k} f \in P_{k}(\Mcal)$ of $f$ provides a suitable approximation satisfying (B5). Furthermore,
\begin{align*}
	E^*(\tau) &= -\int_\Omega W^*(\tau) \d{x} - \chi_{-f} (\div\,\tau) &&\text{for } \tau \in \Sigma,\\
	E^*_h(\tau) &= - \int_\Omega W^*(\tau_\Mcal) \d{x} - \chi_{-f_h}^*(\div_h \tau) - \gamma_h(\tau)/r' &&\text{for } \tau = (\tau_\Mcal, \tau_\mathcal{F}) \in Y
\end{align*}
with the indicator function $\chi_{g}(a) = 0$ if $a = g$ and $+\infty$ if $a \neq g$ for $a,g \in \mathbb{R}$.

\subsection{A~priori}
The ansatz space $V_h$ lacks trace degrees of freedom for the full consistency of the discrete gradient $\nabla_h$ from \eqref{def:discrete-gradient} with respect to discrete test functions.
Therefore, an additional tool is utilized in the a~priori error analysis.
\begin{lemma}[conforming companion]\label{lem:conforming_companion}
	There exists a linear bounded operator $\mathcal{J}_h : V_h \to V$ such that
	any $v_h \in V_h$ satisfies $\Pi_\Mcal^{k} (v_h - J_h v_h) = 0$, $\Pi_S^{k} (J_h v_h - \{v_h\}_S) = 0$ for any $S \in \mathcal{F} \setminus \mathcal{F}_\mathrm{D}$.  Any $K \in \Mcal$ satisfies 
	\begin{align*}
		&h_K^{-1}\|v_h - J_h v_h\|_{L^p(K)}^p\\
		&\qquad + \|\nabla(v_h - J_h v_h)\|_{L^p(K)}^p \lesssim \sum_{S \in \mathcal{F}, S \cap K \neq \emptyset} h^{1-p}\|[v_h]_S\|_{L^p(S)}^p.
	\end{align*}
	In particular, $\nabla_h v_h = \Pi_\Mcal^{k-1} \nabla J_h v_h$.
\end{lemma}
\begin{proof}
	The explicit construction of $J_h$ utilizes well-understood
	averaging and bubble functions techniques, cf.~\cite{VeeserZanotti2018,ErnZanotti2020} for further details. The asserted bound is given in \cite{VeeserZanotti2018,ErnZanotti2020} for $p = 2$ and the general case follows from scaling arguments. Further details on $\mathcal{J}_h$ are omitted.
    
 %    A piecewise integration by parts in \eqref{def:ldg-dicrete-gradient} proves
	% \begin{align*}
	% 	\int_\Omega \nabla_h v_h^\mathrm{ldg} \cdot \Phi \d{x}
	% 	= - \int_\Omega v_h^\mathrm{ldg} \div_\pw \Phi \d{x} + \sum_{S \in \mathcal{F}\setminus\mathcal{F}_\mathrm{D}} \int_S \{v_h^\mathrm{ldg}\}_S [\Phi \cdot \nu_S]_S \d{s}.
	% \end{align*}
	The $L^2$ orthogonality $v_h - J_h v_h \perp \div_\pw \Phi$ and $J_h v_h - \{v_h\}_S \perp [\Phi \cdot \nu_S]_S$ for any $S \in \mathcal{F}\setminus\mathcal{F}_\mathrm{D}$ shows that the right-hand side of \eqref{def:discrete-gradient} is equal to $(\nabla J_h v_h, \Phi)_{L^2(\Omega)}$ for any $\Phi \in P_{k-1}(\Mcal)^n$. This leads to
	$\nabla_h v_h = \Pi_\Mcal^{k-1} \nabla J_h v_h$.
\end{proof}

The subsequent theorem is the main result of this section.

\begin{theorem}[a~priori]\label{thm:error-estimate}
	Suppose (B1), (B4), \eqref{def:lower-order-terms}, and $\sigma \in W^{1,1}(\Omega)^n \cap \Sigma$. Then
	\begin{align*}
		E_h&(\mathrm{I}_h u) - \min E_h(V_h) \leq \int_\Omega (\sigma - \D W(\nabla_h \mathrm{I}_h u)) \cdot (\nabla u - \nabla_h \I_h u) \d{x} + s_h(\I_h u)/r\\
		&+ \int_\Omega \sigma \cdot (\nabla_h \mathrm{I}_h u - \nabla J_h \I_h u) \d{x} + \int_\Omega f (J_h \I_h u - \mathrm{I}_h u) \d{x} + \gamma_h(\mathrm{I}_h^* \sigma)/r'.
	\end{align*}
\end{theorem}

\begin{proof}
\Cref{lem:consistency} implies $\div_h \mathrm{I}_h^* \sigma = -f_h$ and so,
$E^*_h(\mathrm{I}_h^* \sigma) = E^*(\sigma) - \gamma_h(\mathrm{I}_h^* \sigma)/r'$.
This, \Cref{thm:weak-duality}, and \eqref{eq:strong-duality} reveal
	\begin{align}\label{ineq:proof-error-estimate-split}
		0 \leq E_h(\mathrm{I}_h u) - \min E_h(V_h) &\leq E_h(\mathrm{I}_h u) - E^*_h(\mathrm{I}_h^* \sigma)\nonumber\\
		&= E_h(\mathrm{I}_h u) - E(u) + \gamma_h(\mathrm{I}_h^* \sigma)/r'.
	\end{align}
    The convexity $0 \leq W(\nabla u) - W(\nabla_h \I_h u) - \D W(\nabla_h \mathrm{I}_h u) \cdot (\nabla u - \nabla_h \I_h u)$ a.e.~in $\Omega$ of $W$ provides
	\begin{align*}
		E_h(\I_h u) - E(u) \leq - \int_\Omega (\D W(\nabla_h \mathrm{I}_h u) \cdot (\nabla u - \nabla_h \I_h u) + f (u - \I_h u)) \d{x} + \frac{\s_h(\I_h u)}{r}.
	\end{align*}
	The combination of this with the Euler-Lagrange equations
	\begin{align*}
		\int_\Omega f (u - J_h \I_h u) \d{x} = \int_\Omega \sigma \cdot \nabla (u - J_h \I_h u) \d{x}.
	\end{align*}
results in the bound
	\begin{align}\label{ineq:proof-convergence-rates-primal}
		E_h(\I_h u) - E(u) \leq \int_\Omega (\sigma - \D W(\nabla_h \I_h u)) \cdot (\nabla u - \nabla_h \I_h u) \d{x} + \frac{1}{r}\s_h(\I_h u)&\nonumber\\
		+ \int_\Omega \sigma \cdot (\nabla_h \I_h u - \nabla J_h \I_h u) \d{x} + \int_\Omega f (J_h \I_h u - \mathrm{I}_h u) \d{x}&.
	\end{align}
	This and \eqref{ineq:proof-error-estimate-split} conclude the proof.
\end{proof}

Convergence rates in terms of the maximal mesh-size $h_{\max}$ can be derived from \Cref{thm:error-estimate} under suitable smoothness assumptions as follows.

\begin{corollary}[convergence rates]\label{cor:convergence-rates}
Suppose that the assumptions of \Cref{thm:error-estimate} hold and $\D W(\nabla_h \mathrm{I}_h u)$ is uniformly bounded in $L^{p'}(\Omega)^n$ independent of the mesh-size.
	If $u \in V \cap W^{k+1,\max\{p,r\}}(\Mcal)$ and $\sigma \in W^{1,1}(\Omega)^n \cap W^{k,\max\{p',r'\}}(\Mcal)^n$, then $$E_h(\I_h u) - \min E_h(V_h) \lesssim h_{\max}^{\ell}$$
	with $\ell \coloneqq \min\{k, (k+1)r-1-s, ((s+1)+(k-1)r)/(r-1)\}$.
\end{corollary}

\begin{proof}
Standard arguments involving, e.g., the trace inequality and the approximation property of the $L^2$ projections lead to
	\begin{align}\label{ineq:proof-stab-conv}
		\s_h(\mathrm{I}_h u) \lesssim h_\mathrm{max}^{r-1-s}\|\nabla_\pw(u - \Pi_{\Mcal}^k u)\|_r^r &\lesssim h_\mathrm{max}^{(k+1)r-1-s}|u|_{W^{k+1,r}(\Mcal)}^r,\\
		\gamma_h(\mathrm{I}_h^* \sigma) \lesssim h_{\max}^{(s+1)/(r-1)}\|\nabla_\pw(\sigma - \Pi_{\Mcal}^{k-1}\sigma)\|^{r'}_{r'} &\lesssim h_{\max}^{((s+1)+(k-1)r)/(r-1)}|\sigma|_{W^{k,r'}(\Mcal)}^{r'}.\nonumber
	\end{align}
	From \Cref{lem:conforming_companion}, \Cref{lem:discrete-consistency-error}, and a triangle inequality, we deduce that
	\begin{align*}
		h_\mathrm{max}^{-1}\|\mathrm{I}_h u - J_h \mathrm{I}_h u\|_{p} &+ \|\nabla_\pw (\mathrm{I}_h u - J_h \mathrm{I}_h u)\|_{p}\\
		& + \|\nabla_h \I_h u - \nabla J_h \I_h u\|_{p} \lesssim h^k_\mathrm{max}|u|_{W^{k+1,p}(\Mcal)}.
	\end{align*}
	Since $f = -\div\,\sigma \in W^{k-1,p'}(\Mcal)$, this and the $L^2$ orthogonality $\nabla_h \I_h u - \nabla J_h \I_h u \perp P_{k-1}(\Mcal)^n$ and $\I_h u - J_h \mathrm{I}_h u \perp P_{k}(\Mcal)$ from \Cref{lem:conforming_companion} imply
	\begin{align}\label{ineq:proof-conv}
		\int_\Omega \sigma &\cdot (\nabla_h \I_h u - \nabla J_h \I_h u) \d{x} + \int_\Omega f (J_h \I_h u - \mathrm{I}_h u) \d{x}\nonumber\\
		&= \int_\Omega \big((1 - \Pi_\Mcal^{k-1}) \sigma \cdot (\nabla_h \I_h u - \nabla J_h \I_h u) + (1 - \Pi_\Mcal^k) f (J_h \I_h u - \mathrm{I}_h u)\big) \d{x}\nonumber\\
		&\lesssim h^{2k}_\mathrm{max}|u|_{W^{k+1,p}(\Omega)}|\sigma|_{W^{k,p'}(\Omega)}.
	\end{align}
	Under the smoothness assumptions of \Cref{cor:convergence-rates}, \Cref{lem:discrete-consistency-error} leads to $\|\nabla_\pw \mathrm{I}_h u - \nabla_h \mathrm{I}_h u\|_{p} \lesssim h_\mathrm{max}^k|u|_{W^{k+1,p}(\Mcal)}$. This and a triangle inequality prove
	\begin{align*}
		\|\nabla u - \nabla_h \mathrm{I}_h u\|_{p} \lesssim h_\mathrm{max}^k|u|_{W^{k+1,p}(\Mcal)}.
	\end{align*}
	Therefore, a H\"older inequality and the boundedness of $\sigma - \D W(\nabla_h \I_h u)$ in $L^{p'}(\Omega)$ by assumption provide
	\begin{align}\label{ineq:proof-conv-suboptimal}
		\int_\Omega (\sigma &- \D W(\nabla_h \I_h u)) \cdot (\nabla u - \nabla_h \I_h u) \d{x}\nonumber\\
		&\qquad\leq \|\sigma - \D W(\nabla_h \mathrm{I}_h u)\|_{p}\|\nabla u - \nabla_h \mathrm{I}_h u\|_{p}\nonumber\\
		&\qquad\lesssim h_\mathrm{max}^k(\|\sigma\|_{p'} + 1)|u|_{W^{k+1,p}(\Omega)}.
	\end{align}
	The combination of \eqref{ineq:proof-stab-conv}--\eqref{ineq:proof-conv-suboptimal}  with \Cref{thm:error-estimate} concludes the proof.
\end{proof}

\begin{remark}[balancing weights for stabilization]\label{rem:balancing-stab}
	To obtain balanced convergence rates for the stabilizations on the primal and dual level in \eqref{ineq:proof-stab-conv} under the smoothness assumptions of \Cref{cor:convergence-rates}, we can choose the parameter $s = (k+1)(r-2)+1$,
	\begin{align*}
		(k+1)r-1-s = 2k = ((s+1)+(k-1)r)/(r-1).
	\end{align*}
	This leads to quadratic convergence rates for the stabilizations in \eqref{ineq:proof-stab-conv}.
\end{remark}

\begin{remark}[choice of $s$]
Under the assumptions of \Cref{cor:convergence-rates},
the best possible rate is bounded by $k$
obtained for $r-k-2 \leq s \leq (k+1)(r-1)$. 
This includes the choice $r = p$ and $s = p-1$ as in \cite{BurmanErn2008}.
\end{remark}

Under additional structural assumptions on the energy density $W$, however, the convergence rates in \Cref{cor:convergence-rates} can be improved further.
Suppose that $p \geq 2$ and we refer to \Cref{remark:p_smaller_2} below for the case $1 < p < 2$. We consider the assumptions (B1)--(B4) from the introduction.

\begin{remark}[boundedness of primal variable]\label{rem:bdd-primal-variable}
	On the continuous level, the lower growth in (B2) provides the uniform bound
	$\|\nabla u\|_p \lesssim 1$,
	cf.~\cite{CPlechac1997} for explicit constants. This, a triangle inequality, and \Cref{lem:discrete-consistency-error} imply $\|\nabla_h \mathrm{I}_h u\|_p \lesssim 1$.
	Furthermore, \cite[Lemma 2.1(a)]{CarstensenTran2021} provides
	$\|\sigma\|_{p'} + \|\D W(\nabla_h \I_h u)\|_{p'} \lesssim 1$.
\end{remark}

The point is that (B3) implies \cite{GlowinskiMarrocco1975,Chow1989,CPlechac1997},
for all $\alpha,\beta \in L^p(\Omega)^n$, that  
\begin{align}\label{ineq:cc-continuous}
    \delta(\alpha,\beta) &\coloneqq \frac{\|\D W(\alpha) - \D W(\beta)\|_{p'}^2}{(1 + \|\alpha\|^p_{p} + \|\beta\|_{p}^p)^{(2-p')/p'}}\nonumber\\
    &\lesssim \int_\Omega (W(\beta) - W(\alpha) - \D W(\alpha) \cdot (\beta - \alpha)) \d{x}.
\end{align}


\begin{proposition}[convergence rates for degenerate convex minimization problems]\label{prop:conv-rates}
Suppose (B1)--(B4), $s = (k+1)(r-2)+1$, $u \in V \cap W^{k+1,\max\{p,r\}}(\Mcal)$,
    and $\sigma \in W^{1,1}(\Omega)^n \cap W^{k,\max\{p',r'\}}(\Mcal)^n$. Then
	\begin{align*}
		|E(u) - \min E_h(V_h)| + E_h(\mathrm{I}_h u) - \min E_h(V_h) 
        =O(h_\mathrm{max}^{2k}).
%		(|u|_{W^{k+2,p}(\Mcal)} + |\sigma|_{W^{k+1,2}(\Mcal)})
	\end{align*}
\end{proposition}

\begin{proof}
    Exchanging the roles of $\alpha$ and $\beta$ in \eqref{ineq:cc-continuous} followed by the 
    sum of the two resulting inequalities proves, for any $\alpha, \beta \in L^p(\Omega)^n$, that
    \begin{align}\label{ineq:cc-monotonicity}
        \delta(\alpha,\beta) \lesssim \int_\Omega (\D W(\alpha) - \D W(\beta)) \cdot (\alpha - \beta) \d{x}
    \end{align}
	The choice $\alpha \coloneqq \nabla u$ and $\beta \coloneqq \nabla_h \mathrm{I}_h u$ in \eqref{ineq:cc-monotonicity} and a H\"older inequality imply
	\begin{align*}
		\delta(\nabla u, \nabla_h \mathrm{I}_h u) \lesssim \|\sigma - \D W(\nabla_h \mathrm{I}_h u)\|_{p'}\|\nabla u - \nabla_h \mathrm{I}_h u\|_{p}.
	\end{align*} 
Since
$\|\nabla u\|_{L^p(\Omega)} + \|\nabla_h \mathrm{I}_h u\|_{L^p(\Omega)} \lesssim 1$
	from \Cref{rem:bdd-primal-variable}, 
this shows
	\begin{align}\label{ineq:conv-sigma-interp}
		\|\sigma - \D W(\nabla_h \mathrm{I}_h u)\|_{p'} \lesssim \|\nabla u - \nabla_h \mathrm{I}_h u\|_{p} \lesssim h_\mathrm{max}^{k}|u|_{W^{k+1,p}(\Mcal)}.
	\end{align}
Therefore, we deduce from a H\"older inequality that
    \begin{align}\label{ineq:proof-conv-opt}
        \int_\Omega (\sigma - \D W(\nabla_h \I_h u)) \cdot (\nabla u - \nabla_h \I_h u) \d{x} \lesssim h^{2k}|u|_{W^{k+1}(\Mcal)},
    \end{align}
    improving
    the convergence rates in \eqref{ineq:proof-conv-suboptimal}. This, \eqref{ineq:proof-stab-conv}--\eqref{ineq:proof-conv}, \Cref{rem:balancing-stab}, and \Cref{thm:error-estimate} prove
	\begin{align}\label{ineq:proof-conv-rate-1}
		E_h(\mathrm{I}_h u) - \min E_h(V_h) \lesssim h_\mathrm{max}^{2k}.
	\end{align}
It remains to control $|E(u) - E_h(I_h u)|$. Since
$E_h(\mathrm{I}_h u) - E(u) \lesssim h_\mathrm{max}^{2k}$ is known from \eqref{ineq:proof-convergence-rates-primal}--\eqref{ineq:proof-conv} and \eqref{ineq:proof-conv-opt}, it remains to control $E(u) - E_h(\mathrm{I}_h u)$.
The convexity of $W$ implies
$0 \leq W(\nabla_h \mathrm{I}_h u) - W(\nabla u) - \sigma \cdot (\nabla_h \mathrm{I}_h u - \nabla u)$ a.e.~in $\Omega$ and so  
	\begin{align*}
		&E(u) - E_h(\mathrm{I}_h u) \leq -\int_\Omega \sigma \cdot (\nabla_h \mathrm{I}_h u - \nabla u) \d{x} - \int_\Omega f(u - \I_h u) \d{x}\\
		&= - \int_\Omega \sigma \cdot (\nabla_h \mathrm{I}_h u - \nabla J_h \mathrm{I}_h u) \d{x} - \int_\Omega \sigma \cdot (\nabla J_h \mathrm{I}_h u - \nabla u) \d{x} - \int_\Omega f(u - \I_h u) \d{x}.
	\end{align*}
    The identity $(\sigma, \nabla J_h \mathrm{I}_h u - \nabla u)_{L^2(\Omega)} = (f, J_h \mathrm{I}_h u - u)_{L^2(\Omega)}$ from the Euler-Lagrange equations shows that the right-hand side is equal to the negative of the left-hand side of \eqref{ineq:proof-conv}, which implies $E(u) - E_h(\mathrm{I}_h u)
    =O(h_\mathrm{max}^{2k}) $%\lesssim h_\mathrm{max}^{2k}$
    and so, $|E(u) - E_h(\mathrm{I}_h u)| =O(h_\mathrm{max}^{2k}) $.%5\lesssimh_\mathrm{max}^{2k}
    The combination of this with \eqref{ineq:proof-conv-rate-1} and a triangle inequality concludes the proof.
\end{proof}
\begin{remark}[significant choices of $r$]
	For $r = 2$, $s = 1$ in \Cref{prop:conv-rates} is computationally attractive due to its quadratic structure. The case $r = p$ and $s = (k+1)(p-2)$ is of theoretical interest, where the regularity $u \in V \cap W^{k+1,p}(\Mcal)$ and $\sigma \in W^{1,1}(\Omega)^n \cap W^{k,p'}(\Mcal)^n$ is required in \Cref{prop:conv-rates}.
\end{remark}
\begin{remark}[boundedness of discrete primal variable]\label{rem:bound-d-primal-variable}
	Note that the two-sided growth of $W$ in (B2) implies the two-sided growth
	\begin{align}\label{ineq:two-sided-dual}
		c_6|g|^{p'} - c_7 \leq W^*(g) \leq c_8|g|^{p'} + c_9 \quad\text{for any } g \in \R^n.
	\end{align}
	of $W^*$ with positive constants $c_6,c_8 > 0$ and non-negative constants $c_7,c_9 \geq 0$, cf., e.g., \cite[Lemma 2.1(b)]{CarstensenTran2021}. Under the assumptions of \Cref{prop:conv-rates}, the interpolation $\I_h^* \sigma$  satisfies $\div_h \I_h^* \sigma = - f_h$ from \Cref{lem:consistency}. Thus, \eqref{ineq:two-sided-dual} implies
	\begin{align*}
		-c_8\|\sigma\|_{p'}^{p'} - c_9|\Omega| + \gamma_h(\I_h^* \sigma) \leq E_h^*(\I_h^* \sigma) \leq E^*_h(y) \leq -c_6\|\sigma_\Mcal\|_{p'}^{p'} + c_7|\Omega|
	\end{align*}
	This, \eqref{ineq:proof-stab-conv}, and \Cref{rem:balancing-stab} show $\|\sigma_\Mcal\|_{p'} \lesssim 1$. Since $\nabla_h u_h \in \partial W^*(\sigma_\Mcal)$, $\|\nabla_h u_h\|_p \lesssim 1$ from \cite[Lemma 2.1(c)]{CarstensenTran2021}.
\end{remark}
\begin{remark}[convergence rates for the stress error]
	Suppose that the assumptions of \Cref{prop:conv-rates} hold.
	The choice $\alpha \coloneqq \nabla_h u_h$ and $\beta \coloneqq \nabla_h \mathrm{I}_h u$ in \eqref{ineq:cc-continuous} proves
	\begin{align*}
		\delta(\nabla_h u_h,\nabla_h \mathrm{I}_h u) \lesssim \int_\Omega (W(\nabla_h \mathrm{I}_h u) - W(\nabla_h u_h) - \D W(\nabla_h u_h) \cdot \nabla_h(\mathrm{I}_h u - u_h)) \d{x}.
	\end{align*}
	This, the discrete Euler-Lagrange equations \eqref{eq:dELE}, and \Cref{cor:convergence-rates} imply
	\begin{align}\label{ineq:err-quantity-conv}
		 \delta(\nabla_h u_h,\nabla_h \mathrm{I}_h u) \lesssim E_h(\mathrm{I}_h u) - E_h(u_h) \lesssim h_{\mathrm{max}}^{2k}.
	\end{align}
	Since $\nabla_h u_h$ and $\nabla_h \I_h u$ are uniformly bounded in $L^{p}(\Omega)^n$ from \Cref{rem:bdd-primal-variable} and \Cref{rem:bound-d-primal-variable}, $\|\sigma - \D W(\nabla_h u_h)\|_{p'} \lesssim h_{\max}^k$.
\end{remark}
\begin{remark}[strongly monotone]
	Assume (B2)--(B3), and
	\begin{align}
		c_9^{-1}|a - b|^p &\leq W(b) - W(a) - \D W(a) \cdot (b - a) \quad\text{for any } a,b \in \mathbb{R}^n\label{ineq:cc-primal}
		% \begin{split}
			% 	c_2^{-1}|\D W(a) - \D W(b)|^2 &\leq (1 + |a|^{p-2} + |b|^{p-2})\\
			% 	&\quad \times (W(b) - W(a) - \D W(a) \cdot (b - a))
			% \end{split}\label{ineq:cc-dual}
	\end{align}
    and a positive constant $c_9 > 0$.
	Then the abstract error quantity in \eqref{ineq:cc-continuous} can be replaced by
	\begin{align*}
		\delta(\alpha,\beta) \coloneqq \frac{\|\D W(\alpha) - \D W(\beta)\|_{p'}^2}{(1 + \|\alpha\|^p_{p} + \|\beta\|_{p}^p)^{(2-p')/p'}} + \|\alpha - \beta\|_p^p.
	\end{align*}
	Strong convexity of the energy \eqref{def:energy} leads to a unique minimizer $u = \arg\min E(V)$.
	From \eqref{ineq:err-quantity-conv} and a triangle inequality, we deduce the convergence rates $\|\nabla u - \nabla_h u_h\|_p \lesssim h_{\max}^{2k/p}$, improving the rates $h_{\max}^{k/(p-1)}$ over the literature on nonconforming methods of arbitrary order \cite{DiPietroDroniou2017-II,DroniouEymardGallouet2018}.
	The conditions (B2)--(B3) and \eqref{ineq:cc-primal} are satisfied, e.g., in the $p$-Laplace problem.
    Then the convergence rates $h_{\max}^{\min\{2,p'\}}$ for the LDG method of \cite{BurmanErn2008} have been derived in \cite{LDG2014} under regularity assumptions based on the natural distance \cite{EbmeyerLiuSteinhauer2005}. The latter can be guaranteed under natural assumptions on the domain and right-hand side.
    %, while the smoothness assumptions in \Cref{cor:convergence-rates} are unrealistic.
\end{remark} 

%\begin{remark}[two-sided growth]
%	The conditions \eqref{ineq:cc-primal}--\eqref{ineq:cc-dual} imply a two-sided growth of the energy density $W$, that is
%	
%	for any $a \in \mathbb{R}^n$. This result is certainly straight-forward but difficult to find in the literature. Therefore, a short proof is outlined below. By exchanging the roles of $a,b$ in \eqref{ineq:cc-dual}, we infer that
%	\begin{align*}
%		c_2^{-1}|\D W(a) - \D W(b)|^2 \leq (1 + |a|^{p-2} + |b|^{p-2})(\D W(b) - \DW(a)) \cdot (b - a)
%	\end{align*}
%	and so $c_2^{-1}|\D W(a) - \D W(b)| \leq (1 + |a|^{p-2} + |b|^{p-2}) |b - a|$.
%	For $b \coloneqq 0$, a Young and triangle inequality imply
%	\begin{align}\label{ineq:gradient-growth}
%		|\D W(a)| \leq c_7 + c_8|a|^{p-1}.
%	\end{align}
%	with the constants $c_7 = c_2(|\D W(0)| + (p-2)/(p-1))$ and $c_8 = c_2(1 + 1/(p-1))$.
%	The choice $b \coloneqq 0$ in \eqref{ineq:cc-primal} proves $W(a) \leq W(0) + \D W(a) \cdot a$. This, \eqref{ineq:gradient-growth}, and a Young inequality lead to the upper $p$-growth
%	\begin{align*}
%		W(a) \leq c_5|a|^p + c_6
%	\end{align*}
%	with the constants $c_5 = c_8 + 1/p$ and $c_6 = W(0) + c_7^{p'}/p'$.
%	Exchanging the roles of $a$ and $b = 0$ in \eqref{ineq:cc-primal} for convenience implies $c_1^{-1}|a|^p \leq W(a) - W(0) - \D W(0) \cdot a$. This and a Young inequality yield the lower $p$-growth
%	\begin{align*}
%		c_3|a|^p - c_4 \leq W(a)
%	\end{align*}
%	with the constants $c_3 = c_1^{-1}/p'$ and $c_4 = - W(0) + c_1^{p'/p}|\D W(0)|^{p'}/p'$.
%\end{remark}

\begin{remark}[$1 < p \leq 2$]\label{remark:p_smaller_2}
    If $1 < p \leq 2$,
    we assume (B3) with $d(a,b) \coloneqq |\D W(a) - \D W(b)|^{p'}$, which implies \eqref{ineq:cc-monotonicity} for the error quantity
    \begin{align*}
    	\delta(\alpha,\beta) \coloneqq \|\D W(\alpha) - \D W(\beta)\|_{p'}^{p'}.
    \end{align*}
	The choice $\alpha \coloneqq \nabla u$ and $\beta \coloneqq \nabla_h \mathrm{I}_h u$ in \eqref{ineq:cc-monotonicity} and a H\"older inequality lead to
	\begin{align}\label{ineq:conv-sigma-p-sub-2}
		\|\sigma - \D W(\nabla_h \mathrm{I}_h u)\|_{p'} \lesssim \|\nabla u - \nabla \mathrm{I}_h u\|_{p}^{p-1}.
	\end{align}
	Suppose that $u \in V \cap W^{k+1,\max\{p,r\}}(\Mcal)$ and $\sigma \in W^{1,1}(\Omega)^n \cap W^{k,\max\{p',r'\}}(\Mcal)^n$, then \eqref{ineq:conv-sigma-p-sub-2} and a H\"older inequality imply
	\begin{align*}
		\int_\Omega (\sigma - \D W(\nabla_h \I_h u)) \cdot (\nabla u - \nabla_h \I_h u) \d{x} \leq \|\nabla u - \nabla_h \I_h u\|_p^p \lesssim h_{\max}^{kp}.
	\end{align*}
	The combination of this with \eqref{ineq:proof-stab-conv}--\eqref{ineq:proof-conv}, and \Cref{thm:error-estimate} concludes
	\begin{align}\label{ineq:conv-rate-p-sub-2}
		|E(u) - \min E_h(V_h)| + E_h(\mathrm{I}_h u) - \min E_h(V_h) \lesssim h_{\max}^{kp}
	\end{align}
	for $kp(r-1) - (k-1)r - 1 \leq s \leq (k+1)r - kp - 1$.
	This, $\delta(\nabla_h u_h, \nabla_h \I_h u) \lesssim E_h(\I_h u) - E_h(u_h)$ from \eqref{ineq:err-quantity-conv}, \eqref{ineq:conv-sigma-p-sub-2}, and a triangle inequality conclude
	\begin{align*}
		\|\sigma - \D W(\nabla_h u_h)\|_{p'} \lesssim h_{\max}^{k(p-1)}.
	\end{align*}
	If we additionally assume that any $a,b \in \mathbb{R}^n$ satisfy
	\begin{align*}
		\frac{|a-b|^2}{1 + |a|^{2-p} + |b|^{2-p}} \lesssim W(b) - W(a) - \D W(a) \cdot (b-a),
	\end{align*}
	then any $\alpha, \beta \in L^p(\Omega)^n$ satisfy \eqref{ineq:cc-continuous} with
	\begin{align*}
		\delta(\alpha,\beta) \coloneqq \|\D W(\alpha) - \D W(\beta)\|_p^p + \frac{\|\alpha - \beta\|^2}{(1 + \|\alpha\|_p^p + \|\beta\|_p^p)^{(2-p)/p}}.
	\end{align*}
	From \eqref{ineq:err-quantity-conv} and \eqref{ineq:conv-rate-p-sub-2}, we infer
    $\delta(\nabla_h u_h,\nabla_h \mathrm{I}_h u) \lesssim E_h(\mathrm{I}_h u) - E_h(u_h) = O(h_{\max}^{kp})$.
	The uniform boundedness of $\|\nabla_h u_h\|_{p} \lesssim 1$ follows from the arguments of \Cref{rem:bound-d-primal-variable}. This, $\delta(\nabla_h u_h,\nabla_h \mathrm{I}_h u) \lesssim h_{\max}^{kp}$, and a triangle inequality imply the convergence rates $\|\nabla u - \nabla_h u_h\|_p \lesssim h_\mathrm{max}^{kp/2}$. This recovers the result of \cite{Tran2024} for hybridizable and \cite{Chow1989} for conforming methods for strongly monotone problems.
\end{remark}

\subsection{A~posteriori}\label{sec:a-posteriori}
A $\Sigma$-conforming approximation of the dual variable in the Raviart-Thomas finite element space is constructed by direct prescription of the degrees of freedom. This provides an alternative to equilibrium techniques \cite{LuceWohlmuth2004,BraessSchoeberl2008,ErnVohralik2015} with solving local problems.
For the sake of brevity, we assume that $\Mcal$ is a regular triangulation into simplices (without hanging nodes) and refer to \cite[Section 5]{Tran2024} for further details if $\Mcal$ is a polytopal mesh.

Recall $y = (\sigma_\Mcal, \sigma_\mathcal{F}) \in Y$ from \eqref{def:dual-variable}. 
Let $\sigma_\RT \in \mathrm{RT}_{k}(\Mcal) \cap \Sigma$ be the unique Raviart-Thomas finite element function with
\begin{align}\label{def:sigma-h}
	\Pi_\Mcal^{k-1} \sigma_\RT = \Pi_\Mcal^{k-1} \sigma_\Mcal \quad\text{and}\quad \Pi_S^k \sigma_\RT = \Pi_S^k \sigma_\mathcal{F} \text{ for any } S \in \mathcal{F}\setminus\mathcal{F}_\mathrm{N}.
\end{align}

\begin{theorem}[post-processing]\label{thm:post-processing}
	Let $\Mcal$ be a regular triangulation of $\Omega$ into simplices.
	Suppose (B1) and $\psi_h(x,\bullet) \in C^1(\R)$ for a.e.~$x \in \Omega$. Then $\sigma_\RT \in \RT_k(\Mcal) \cap \Sigma$ from \eqref{def:sigma-h} satisfies 
    $\div\,\sigma_\RT = \div_h \sigma_h = \Pi_\Mcal^{k} \partial_u \psi_h(\bullet,u_h).$
\end{theorem}

\begin{proof}
	% The discrete Euler-Lagrange equations associated with the minimization problem \eqref{def:energy-ldg} and a straightforward modification of \eqref{eq:diby-hho} imply
	% \begin{align}\label{eq:ldg-ELE}
	% 	0 &= \int_\Omega \sigma_\Mcal^\mathrm{ldg} \cdot \nabla_h^\mathrm{ldg} v_h^\mathrm{ldg} \d{x} + \int_\Omega \nabla_u \psi_h(x, u_h^\mathrm{ldg}) v_h^\mathrm{ldg} \d{x} + s_h^\mathrm{ldg}(u_h^\mathrm{ldg};v_h^\mathrm{ldg})\nonumber\\
	% 	&= - \int_\Omega (\nabla_u \psi_h(x,u_h^\mathrm{ldg}) - \div_h \sigma_h^\mathrm{ldg}) v_h^\mathrm{ldg} \d{x}\nonumber\\
	% 	&\qquad + \sum_{S \in \mathcal{F} \setminus \mathcal{F}_\mathrm{N}} \int_S [v_h^\mathrm{ldg}]_S (\sigma_S^\mathrm{ldg} - \{\Pi_\Mcal^k \sigma_\Mcal^\mathrm{ldg}\}_S \cdot \nu_S) \d{s} + s_h^\mathrm{ldg}(u^\mathrm{ldg}_h; v_h^\mathrm{ldg})
	% \end{align}
	% for any $v_h^\mathrm{ldg} \in V_h^\mathrm{ldg}$.
	% By definition of $\sigma_\mathcal{F}^\mathrm{ldg}$, the sum of the final two terms on the right-hand side of \eqref{eq:ldg-ELE} vanishes and so, \eqref{eq:ldg-ELE} implies the first assertion $\div_h \sigma_h^\mathrm{ldg} = \Pi_\Mcal^{k+1} \nabla_u \psi(x,u_h^\mathrm{ldg})$.
Integration by parts, \eqref{def:sigma-h}, and \eqref{def:div-rec} show, for any $\phi \in P_k(\Mcal)$, that
	\begin{align*}
		\int_\Omega \div\,\sigma_\RT \phi \d{s} &= -\int_\Omega \sigma_\RT \cdot \nabla_\pw \phi \d{x} + \sum_{S \in \mathcal{F}\setminus\mathcal{F}_\mathrm{N}} \int_S [\phi]_S \sigma_\RT\cdot \nu_S \d{s}\\
		&= -\int_\Omega \sigma_\Mcal \cdot \nabla_\pw \phi \d{x} + \sum_{S \in \mathcal{F}\setminus\mathcal{F}_\mathrm{N}} \int_S [\phi]_S \sigma_S \d{s} = \int_\Omega \div_h \sigma_h \phi \d{x}.
	\end{align*}
	Since $\div_h \sigma_h = \Pi_\Mcal^{k} \partial_u \psi_h(\bullet,u_h)$ by \eqref{eq:div_h_sigma_h}, this concludes the proof.
\end{proof}
% \begin{remark}[strong duality to dual problem]
% 	Suppose that the assumptions of \Cref{thm:strong-duality} hold.
% 	The arguments of \Cref{sec:equivalence} and the computation in the proof of \Cref{thm:post-processing} show that the LDG method of \cite{BurmanErn2008} is equivalent to the dual problem that maximizes
% 	\begin{align*}
% 		\widetilde{E}^*_h(\tau_h) \coloneqq -\int_\Omega (W^*(\tau_\Mcal) + \psi^*_h(x,\div_h \tau_h)) - \widetilde{\gamma}_h(\tau_h)
% 	\end{align*}
% 	among $\tau_h = (\tau_\Mcal, \tau_\mathcal{F}) \in W$, where
% 	\begin{align*}
% 		\widetilde{\gamma}_h(\tau_h) \coloneqq \sum_{K \in \Mcal} \sum_{S \in \mathcal{F}(K) \setminus \mathcal{F}_\mathrm{N}} h^{s/(r-1)} \|\tau_S - \Pi_K^k \tau_\Mcal\|_{L^{r'}(S)}^{r'}
% 	\end{align*}
% 	is a modified stabilization on the dual level without penalization on Neumann boundary sides. Obviously, $\max E^*_h(W) \leq \max \widetilde{E}^*_h(W)$ but it is not known whether the reverse direction holds for nonempty Neumann boundary (which would imply $\min E_h(V_h) = \min E_h^\mathrm{ldg}(V_h^\mathrm{ldg})$).
% \end{remark}
\begin{remark}[a~posteriori error control]
Suppose that the lower-order term $\psi$ has the explicit representation (B4) and the assumptions of \Cref{thm:post-processing} hold.
	Furthermore, we assume that $f \in P_k(\Mcal)$ is a piecewise polynomial.
	Given a conforming postprocessing $v_C \in V$, then the energy error $E(v_C) - E(u)$ can be bounded by
	\begin{align*}
		E(v_C) - E(u) \leq E(v_C) - E^*(\sigma_\RT).
	\end{align*}
In the numerical examples below, $v_C$ is obtained from the discrete minimizer $u_h$ of $E_h$ in $V_h$ by nodal averaging in the conforming subspace $V_h \cap V$. If the energy density $W$ satisfies further structural properties, e.g., \eqref{ineq:cc-continuous}, then the Euler-Lagrange equations and the previously displayed formula imply
	\begin{align}\label{ineq:a-posteriori}
		\begin{split}
			\delta(\nabla u, \nabla v_C) &\lesssim \int_\Omega (W(\nabla v_C) - W(\nabla u) - \sigma \cdot \nabla (v_C - u)) \d{x}\\
			&= E(v_C) - E(u) \leq E(v_C) - E^*(\sigma_\RT) \eqqcolon \eta.
		\end{split}
	\end{align}
If $f$ is not piecewise polynomial, then additional data oscillation arises in \eqref{ineq:a-posteriori}. However, this additional error term can be computed explicitly \cite[Remark 5.3]{Tran2024}.
\end{remark}

\section{Extension to hybridizable method}\label{sec:hybrid}
In this section, we briefly extend the analysis of \Cref{sec:error-analysis} to a hybridizable method using the techniques of \cite{Tran2024}. For the sake of simplicity, we retain the notation of \Cref{sec:discretization} on the discrete level.
Given $k \geq 1$, let
\begin{align*}
	V_h \coloneqq P_k(\Mcal) \times P_k(\mathcal{F}\setminus\mathcal{F}_\mathrm{D})
\end{align*}
denote the discrete ansatz space. Given $v_h = (v_\Mcal,v_\mathcal{F}) \in V_h$, the discrete gradient $\nabla_h v_h \in P_{k-1}(\Mcal)^n$ of $v_h$ is the unique solution to
\begin{align*}
	\int_\Omega \nabla_h v_h \cdot \Phi \d{x} &= -\int_\Omega v_\Mcal \,\div_\pw \Phi \d{x} + \sum_{S \in \mathcal{F}\setminus\mathcal{F}_\mathrm{D}} \int_S v_S [\Phi \cdot \nu_S] \d{s}
\end{align*}
for any $\Phi \in P_{k-1}(\Mcal)$,
where $v_S = v_\mathcal{F}|_S$ abbreviates the restriction of $v_\mathcal{F}$ along the side $S$. The discrete problem minimizes
\begin{align}
	E_h(v_h) \coloneqq \int_\Omega (W(\nabla_h v_h) + \psi_h(\bullet,v_\Mcal)) \d{x} + \s_h(v_h)/r
\end{align}
among $v_h = (v_\Mcal, v_\mathcal{F}) \in V_h$
with the stabilization
\begin{align*}
	\s_h(v_h) \coloneqq \sum_{K \in \Mcal} \sum_{S \in \mathcal{F}(K)} h_S^{-s} \int_S T_{K,S} v_h (v_S - v_K) \d{s}
\end{align*}
and $T_{K,S} v_h \coloneqq |v_S - v_K|^{r-2}(v_S - v_K)$ for any $K \in \Mcal$, $S \in \mathcal{F}$.
The corresponding dual problem is \eqref{def:dual-discrete-energy}, 
but with the stabilization
\begin{align*}
	\gamma_h(\tau) \coloneqq \sum_{K \in \Mcal} \sum_{S \in \mathcal{F}} h^{s/(r-1)}\|\tau_S - \Pi_\Mcal^{k-1}\tau_K \cdot \nu_S\|^{r'}_{L^{r'}(S)}
\end{align*}
for any $\tau = (\tau_\Mcal, \tau_\mathcal{F}) \in Y$
instead of \eqref{def:dual-stab}
to reflect the hydridization of the ansatz space.
Here, $\tau_K = \tau_\Mcal|_K$ is the restriction of $\tau_\Mcal$ to $K$.
\begin{remark}[suboptimal polynomial consistency]\label{rem:LS-stab}
	For the Lehrenfeld-Sch\"oberl stabilization, we can use the discrete ansatz space $P_k(\Mcal) \times P_{k-1}(\mathcal{F}\setminus\mathcal{F}_\mathrm{D})$, reducing the computational cost of the method. However, the analysis of this section does not carry over because \eqref{def:div-rec} forfeits to hold.
\end{remark}

The following  \Cref{thm:duality-hybrid} allows for the extension of {\em all results of \Cref{sec:error-analysis}} to the hybrid method this section.


\begin{theorem}[duality of hybridizable methods]\label{thm:duality-hybrid}
	It holds $\sup E^*_h(Y) \leq \min E_h(V_h)$; 
	(B1) and (B5) imply
	%If $\ell = k$, we additional assume that $r = 2$.
    $\max E_h^*(Y) = \min E_h(V_h)$.
\end{theorem}

\begin{proof}
For any $v_h = (v_\Mcal, v_\mathcal{F}) \in V_h$ and $\tau = (\tau_\Mcal, \tau_\mathcal{F}) \in Y$, the proof departs from the integration by parts formula
	\begin{align}\label{eq:dip-hho}
		\int_\Omega \tau_\Mcal \cdot \nabla_h v_h \d{x} &= \int_\Omega \Pi_{\Mcal}^{k-1} \tau_\Mcal \cdot \nabla_h v_h \d{x} = - \int_\Omega v_\Mcal \cdot \div_h \tau \d{x}\nonumber\\ 
		&\qquad+ \sum_{S \in \mathcal{F}\setminus\mathcal{F}_\mathrm{D}} \int_S (v_S - \{v_\Mcal\}_S) \cdot [\Pi_{\Mcal}^{k-1} \tau_\Mcal \cdot \nu_S]_S \d{s}\nonumber\\
		&\qquad + \sum_{S \in \mathcal{F}\setminus\mathcal{F}_\mathrm{N}} \int_S [v_\Mcal]_S \cdot (\tau_S - \{\Pi_{\Mcal}^{k-1} \tau_\Mcal \cdot \nu_S\}_S) \d{s}
	\end{align}
This follows from arguments similar to \cite[Lemma 3.2]{Tran2024}. 
Furthermore, the final two sums on the right-hand side can be rewritten as
	\begin{align*}
		- \sum_{K \in \mathcal{M}} \sum_{S \in \mathcal{F}(K)} (\nu_S \cdot \nu_K) \int_S (v_S - v_K) \cdot (\tau_S - \Pi_\Mcal^{k-1}\tau_K \cdot \nu_S) \d{s},
	\end{align*}
	cf.~\cite[Proof of Theorem 3.1]{Tran2024}. This, \eqref{eq:dip-hho}, the H\"older inequality, and $\tau_\Mcal \cdot \nabla_h v_h \leq W(\nabla_h v_h) + W^*(\tau_\Mcal)$ as well as $\div_h \tau \, v_h \leq \psi_h(\bullet,v_h) + \psi_h^*(\bullet,\div_h \tau_h)$ a.e.~in $\Omega$ conclude $\sup E^*_h(Y) \leq \min E_h(V_h)$.
    
	Assuming the differentiability of $W$ and $\psi_h$, we can define the stress variable $y = (\sigma_\Mcal, \sigma_\mathcal{F}) \in Y$ as
	\begin{align}\label{def:dual-variable-hybrid}
		\begin{split}
			\sigma_\Mcal &\coloneqq \D W(\nabla_h u_h),\\
			\sigma_\mathcal{F}|_S &\coloneqq \begin{cases}
				\{\Pi_\Mcal^{k-1} \sigma_\Mcal\}_S \cdot \nu_S + h_S^{-s}(T_{K_+,S} u_h - T_{K_-,S} u_h)/2 &\mbox{if } S \in \mathcal{F}(\Omega),\\
				\Pi_\Mcal^{k-1} \sigma_\Mcal \cdot \nu_S + h_S^{-s}T_{K_+,S} u_h &\mbox{if } S \in \mathcal{F}_\mathrm{D}.
			\end{cases}
		\end{split} 
	\end{align}
The computations in \cite[Corollary 5.1 and Lemma 5.2]{Tran2024} carry over and show that
	\begin{align}
		[\sigma_\Mcal \nu_S]_S &= - \sum_{K \in \Mcal, S \in \mathcal{F}(K)} h_S^{-s} T_{K,S} u_h \quad\text{for any } S\in\mathcal{F}\setminus\mathcal{F}_\mathrm{N},\label{eq:jump-dual-var}\\
		\div_h y &= \partial_u \psi_h(\bullet,u_\Mcal).\nonumber
	\end{align}
	Thus, the discrete Euler-Lagrange equations and \eqref{eq:dip-hho} imply
	\begin{align}\label{eq:proof-strong-dual-hybrid}
		- \int_\Omega \div_h y\, u_\Mcal \d{x} & = - \int_\Omega \partial_u \psi_h(\bullet,u_\Mcal) u_\Mcal \d{x}\nonumber\\
		& = \int_\Omega \sigma_\Mcal \cdot \nabla_h u_h \d{x} + \s_h(u_h).
	\end{align}
	An explicit computation with the definition of $y$ in \eqref{def:dual-variable-hybrid} and \eqref{eq:jump-dual-var} show, for any $K \in \Mcal$ and $S \in \mathcal{F}(K)$, that
	\begin{align*}
		\sigma_S - \Pi_K^{k-1} \sigma_K \cdot \nu_S &= -(\nu_S \cdot \nu_K)[\Pi_\Mcal^{k-1} \sigma_\Mcal]_S/2 + h_S^{-s}(T_{K_+,S} u_h - T_{K_-,S} u_h)/2\\
        &= (\nu_S \cdot \nu_K)h^{-s} T_{K,S} u_h.
	\end{align*}
	for interior sides $S \in \mathcal{F}(\Omega)$ and
	\begin{align*}
		\sigma_S - \Pi_K^{k-1} \sigma_K \cdot \nu_S = h_S^{-s} T_{K,S} u_h
	\end{align*}
	for boundary sides $S \in \mathcal{F}(\partial \Omega)$.
	This and $s/(r-1) - sr' = -s$ imply
	\begin{align}\label{eq:stab-hybrid-strong-dual}
		\gamma_h(y) = \s_h(u_h).
	\end{align}
	From \eqref{eq:proof-strong-dual-hybrid}--\eqref{eq:stab-hybrid-strong-dual}, $\sigma_\Mcal \cdot \nabla_h u_h = W(\nabla_h u_h) + W^*(\sigma_\Mcal)$, and $\div_h y \,u_\Mcal = \psi_h(\bullet,u_\Mcal) + \psi_h^*(\bullet,\div_h y)$ a.e.~in $\Omega$, we conclude $E_h^*(y) = E_h(u_h)$ and thus, $\max E_h^*(Y) = \min E_h(V_h)$.
\end{proof}

\section{Numerical examples}\label{sec:numerical_examples}
This section tests the performance of the a~posteriori error 
control \eqref{ineq:a-posteriori} in three numerical benchmarks 
in the L-shaped domain $\Omega \coloneqq (-1,1)^2 \setminus ([0,1) \times (-1,0])$ with constant right-hand side $f \equiv 1$.
The initial triangulation in all benchmarks is displayed in 
\Cref{fig:volume_fraction}(a).
The computer experiments are carried out on regular triangulations into \emph{simplices}.

\subsection{Adaptive mesh-refining algorithm}
Since $f$ is constant, the a~posteriori error estimator \eqref{ineq:a-posteriori} applies without data oscillation. The following localization of the right-hand side of \eqref{ineq:a-posteriori} was discussed in \cite{BartelsKaltenbach2023}.
An integration by parts with $\div \sigma_\RT = - f$ implies
\begin{align*}
	\eta = E(v_C) - E^*(\sigma_\RT) = \int_\Omega (W(\nabla v_C) - \sigma_\RT \cdot \nabla v_C + W^*(\sigma_\RT)) \d{x} \eqqcolon \sum_{K \in \Mcal} \eta(K)
\end{align*}
with the local refinement indicator
\begin{align}\label{def:refinement-indicator}
	\eta(K) \coloneqq \int_K (W(\nabla v_C) - \sigma_\RT \cdot \nabla v_C + W^*(\sigma_\RT)) \d{x}.
\end{align}
The Fenchel-Young inequality $W(\nabla v_C) - \sigma_\RT \cdot \nabla v_C + W^*(\sigma_\RT) \geq 0$ holds pointwise a.e.~in $\Omega$, whence $\eta(K) \geq 0$.
Adaptive computations utilize the refinement indicator \eqref{def:refinement-indicator} in the standard adaptive mesh-refining loop \cite{Doerfler1996,CarstensenFeischlPagePraetorius2014} with the D\"orfler marking strategy, i.e., at each refinement step, a subset $\mathfrak{M} \subset \mathcal{M}$ with minimal cardinality is selected such that
\begin{align*}
	\sum\nolimits_{K \in \mathcal{M}} \eta(K) \leq \frac{1}{2}\sum\nolimits_{K \in \mathfrak{M}} \eta(K).
\end{align*}
The convergence history plots display the a~posteriori error estimator $E(v_C) - E^*(\sigma_\RT)$ against the number of degrees of freedom $\mathrm{ndof}$ in a log-log plot.
(Recall the scaling $\mathrm{ndof} \approx h^{-2}_\mathrm{max}$ for uniform meshes.)
Solid lines indicate adaptive, while dashed lines are associated with uniform mesh refinements.
All plotted adaptive meshes are generated with the polynomial degree $k = 2$.

The discrete minimization problem $\min E_h(V_h)$ from \eqref{def:discrete_energy} is solved by an iterative solver \texttt{fminunc} from the MATLAB standard library in an extension of the data structures and the short MATLAB programs \cite{AlbertyCFunken1999}.
The first and (piecewise) second derivatives of $W$ have been provided for 
the trust-region quasi-Newton scheme with \texttt{MaxIterations} = $10^3$, while
\texttt{FunctionTolerance},  
\texttt{OptimalityTolerance}, and 
\texttt{tepTolerance}
in \texttt{fminunc} are set to 
$10^{-15}$.
The numerical integration of piecewise polynomials is carried out exactly.

For non-polynomial functions such as $W(\nabla_h v_h)$ with $v_h \in V_h$, the number of chosen quadrature points allows for exact integration of polynomials of degree at most $2pk + 1$ with the growth $p$ of $W$ and the polynomial order $k$ of the discretization. On the initial triangulation, the starting point for $\texttt{fminunc}$ is zero, while the conforming postprocessing $v_C$ 
initializes the starting point for the finer mesh. 

\begin{figure}[ht]
	\begin{minipage}[b]{0.475\textwidth}
		\centering
		\includegraphics[height=4.6cm]{figures/initial_triangulation}
	\end{minipage}\hfill
	\begin{minipage}[b]{0.475\textwidth}
		\centering
		\includegraphics[height=4.6cm]{figures/odp/material_distribution}
	\end{minipage}
	\captionsetup{width=1\linewidth}
	\caption{(a) Initial triangulation of the L-shaped domain into 6 triangles and (b) material distribution in the optimal design problem of \Cref{sec:odp}}
	\label{fig:volume_fraction}
\end{figure}
\begin{figure}[ht]
	\begin{minipage}[b]{0.475\textwidth}
		\centering
		\includegraphics[height=5cm]{figures/odp/convergence_odp_a_posteriori}
	\end{minipage}\hfill
	\begin{minipage}[b]{0.475\textwidth}
		\centering
		\includegraphics[height=5cm]{figures/odp/adapt_tri}
	\end{minipage}
	\captionsetup{width=1\linewidth}
	\caption{(a) Convergence history plot of $\eta$ for $k = 1, \dots, 4$ and (b) adaptive triangulation in the optimal design problem of \Cref{sec:odp}}
	\label{fig:odp-conv}
\end{figure}

\subsection{Optimal design problem}\label{sec:odp}
This model problem seeks the optimal distribution of two materials with fixed amounts to fill a given domain for maximal torsion stiffness \cite{KohnStrang1986,BartelsC2008}. Given parameters $0 < t_1 < t_2$ and $0 < \mu_1 < \mu_2$ with $t_1 \mu_2 = \mu_1 t_2$, the energy density $W(a) \coloneqq w(|a|)$, $a \in \mathbb{R}^2$, with
\begin{align*}
	w(t) \coloneqq \begin{cases}
		\mu_2 t^2/2 &\mbox{if } 0 \leq t \leq t_1,\\
		t_1\mu_2(t - t_1/2) &\mbox{if } t_1 \leq t \leq t_2,\\
		\mu_1 t^2/2 + t_1\mu_2(t_2/2 - t_1/2) &\mbox{if } t_2 \leq t
	\end{cases}
\end{align*}
satisfies \eqref{ineq:a-posteriori} with $\delta(\alpha,\beta) \coloneqq \|\D W(\alpha) - \D W(\beta)\|^2_2$ for any $\alpha,\beta \in L^2(\Omega)^2$. Therefore, the energy error $E(v_C) - E^*(\sigma_\RT)$ provides an upper bound for the stress error $\|\sigma - \D W(\nabla v_C)\|^2_2$.
This benchmark considers the parameters $\mu_1 = 1$, $\mu_2 = 2$, $t_1 = \sqrt{2\lambda\mu_1/\mu_2}$ for $\lambda = 0.0145$, $t_2 = \mu_2 t_1/\mu_1$ from \cite{BartelsC2008}, the input $r = 2$, $s = 1$ for the stabilization $\s_h$, and $\Gamma_\mathrm{D} = \partial \Omega$.

The approximated material distribution in the  adaptive computation with $k = 1$ is displayed in \Cref{fig:volume_fraction}(b) using  volume fraction plot
\cite[Section 5]{BartelsC2008}.
On uniform meshes, a convergence rate  $2/3$ for $\eta$ is observed 
in \Cref{fig:odp-conv}(a).
The adaptive algorithm refines towards the singularity at the origin and the transition layer in \Cref{fig:odp-conv}(b).
This leads to improved convergence rates for $\eta$ although the improvements appear marginal for higher polynomial degrees.

\begin{figure}[ht]
	\begin{minipage}[b]{0.475\textwidth}
		\centering
		\includegraphics[height=5cm]{figures/plaplace/convergence_plaplace_err}
	\end{minipage}\hfill
	\begin{minipage}[b]{0.475\textwidth}
		\centering
		\includegraphics[height=5cm]{figures/plaplace/adapt_tri_2058_triangles}
	\end{minipage}
	\captionsetup{width=1\linewidth}
	\caption{(a) Convergence history plot of $\eta$ for $k = 1, \dots, 4$ and (b) adaptive triangulation in the $4$-Laplace problem of \Cref{sec:plaplace}}
	\label{fig:plaplace_conv}
\end{figure}
\subsection{$4$-Laplace problem}\label{sec:plaplace}
In this benchmark, we consider the $4$-Laplace problem with $W(a) \coloneqq |a|^4/4$
for any $a \in \mathbb{R}^2$ in $\Omega$ and Dirichlet boundary $\Gamma_\mathrm{D} \coloneqq (\{0\} \times [-1,0] \cup [0,1] \times \{0\})$.
Since $W$ satisfies (B3) and \eqref{ineq:cc-primal}, \eqref{ineq:a-posteriori} holds with
\begin{align*}
	\delta(\nabla u, \nabla v_C) \coloneqq \|\nabla(u - v_C)\|^4_4 + \frac{\|\D W(\nabla u) - \D W(\nabla v_C)\|^2_{4/3}}{(1 + \|\nabla u\|_4^4 + \|\nabla v_C\|_4^4)^{1/2}}.
\end{align*}
Furthermore, $\D W$ is strongly monotone w.r.t.~the quasi norm \cite{BarrettLiu1993,BarrettLiu1994,DieningKreuzer2008} so that, additionally, the error $\|\sqrt{\varrho}\nabla(u - v_C)\|^2_2$ with $\varrho \coloneqq (|\nabla u| + |\nabla v_C|)^{2}$ can be controlled.

On uniform meshes, \Cref{fig:plaplace_conv}(a) displays the convergence rates $2/3$ for $\eta$. Adaptive computation refines towards the re-entrant corner  in \Cref{fig:plaplace_conv}(b) recover the optimal convergence rates $\mathrm{ndof}^{-k}$ for all displayed polynomial degrees $k$. For $k = 1$, the empirical results are consistent with the known optimality of adaptive algorithms for P1 conforming discretizations in \cite{DieningKreuzer2008,BelenkiDieningKreuzer2012}.

\begin{figure}[ht]
	\begin{minipage}[b]{0.475\textwidth}
		\centering
		\includegraphics[height=5cm]{figures/bingham/convergence_bingham_err}
	\end{minipage}\hfill
	\begin{minipage}[b]{0.475\textwidth}
		\centering
		\includegraphics[height=5cm]{figures/bingham/adap_tri_2112.jpg}
	\end{minipage}
	\captionsetup{width=1\linewidth}
	\caption{(a) Convergence history plot of $\eta$ for $k = 1, \dots, 4$ and adaptive triangulation in the Bingham flow problem of \Cref{sec:bingham}. The results are obtained with $\varepsilon = 10^{-5}$}
	\label{fig:bingham}
\end{figure}
\begin{figure}[ht]
	\includegraphics[height=8cm]{figures/bingham/convergence_bingham_err_reg}
	\captionsetup{width=0.8\linewidth}
	\caption{Convergence history plot of $\eta$ for $k = 2$ and $\varepsilon = 10^{-3}, \dots, 10^{-6}$ in \Cref{sec:bingham}}
	\label{fig:bingham-II}
\end{figure}


\subsection{Bingham flow through a pipe}\label{sec:bingham}
Given fixed positive parameters $\mu, g > 0$, the modelling of a uni-directional flow through a pipe with cross-section $\Omega \subset \mathbb{R}^2$ leads to the minimization problem \eqref{def:energy} with the energy density
\begin{align*}
	W(a) \coloneqq \mu|a|^2/2 + g|a| \quad\text{for any } a \in \mathbb{R}^2,
\end{align*} 
cf.~\cite{DuvantLions1972,CarstenReddySchedensack2016}, and $\Gamma_\mathrm{D} = \partial \Omega$.
An explicit computation \cite{Tran2024} shows that
\begin{align*}
	W^*(\alpha) = \begin{cases}
		0 &\mbox{if } |\alpha| \leq g\\
		(|\alpha| - g)^2/(2 \mu) &\mbox{if } |\alpha| > g.
	\end{cases}
\end{align*}
The strict convexity of $W$ leads to a unique the minimizer $u$ of $E$ in $V$. Although $W$ is not differentiable, there exists $\sigma \in H(\div,\Omega) = W^2(\div,\Omega)$ such that $\sigma \in \partial W(\nabla u)$ and $\div\, \sigma = -f$ pointwise a.e.~in $\Omega$ \cite[Chapter II, Theorem 6.3]{Glowinski2008}.
Thus, there is no duality gap $E(u) = E^*(\sigma)$.
Furthermore, \eqref{ineq:a-posteriori} is satisfied with $\delta(\alpha,\beta) \coloneqq \|\alpha - \beta\|^2$ for any $\alpha, \beta \in L^2(\Omega)^2$ \cite[Lemma 1]{CarstenReddySchedensack2016}.

The postprocessings for \eqref{ineq:a-posteriori} are obtained from a regularized discrete problem as in \cite{CarstenReddySchedensack2016,Tran2024}. Given $\varepsilon > 0$, define $W_\varepsilon \in C^1(\mathbb{R}^2)$ by
\begin{align*}
	W_\varepsilon(a) \coloneqq \mu|a|^2/2 + g(\sqrt{|a|^2 + \varepsilon^2}) \quad\text{for any } a \in \mathbb{R}^2.
\end{align*}
The unique minimizer $u_{h,\varepsilon} \in V_h$ of the discrete energy
\begin{align*}
	E_{h,\varepsilon}(v_h) \coloneqq \int_\Omega (W_\varepsilon(\nabla_h v_h) - f v_h) \d{x} + \s_h(v_h)/2
\end{align*}
among $v_h \in V_h$
allows for the postprocessings $\sigma_\RT \in H(\div,\Omega)$ with $\div\,\sigma_\RT = -f$ from as in \Cref{thm:post-processing} and $v_C \in V_h \cap V$ as the nodal average of $u_{h,\varepsilon}$.

The computer experiment runs $\mu = 1$, $g = 0.2$, $f \equiv 1$, and 
shows the convergence rate 0.8  on uniform meshes in \Cref{fig:bingham}(a). 
Adaptive computation refines towards a parameter dependent region and towards 
the re-entrant corner. This leads to a significant improvement for higher-order discretizations $k \geq 2$. Empirical convergence rates are difficult to determine as a plateau is reached due to the regularization. \Cref{fig:bingham}(b) displays the effect of regularization on the error estimator for different 
parameters $\varepsilon$.

\subsection{Conclusions}
The numerical experiments of this section provide similar empirical results to those of \cite{Tran2024} with the hybridizable method outlined in \Cref{rem:LS-stab}. Adaptive mesh-refining leads to improved convergence rates for the a~posteriori error estimator compared to uniform mesh refinements. For the $p$-Laplace problem, optimal convergence rates are recovered with a quadratic stabilization by adaptive mesh-refining algorithms for singular solutions. In fact we
recommend $r=2$ for nonlinear problems. 

%% bibliography
\printbibliography

\appendix

\section{A~priori error analysis of conforming methods}\label{appendix}
Energy error of conforming methods for the convex minimization \eqref{def:energy}
are certainly understood from the arguments of \cite{GlowinskiMarrocco1975,Chow1989,CPlechac1997}, but precise statements are 
rare and provided for completeness in this appendix to explain  \eqref{confenergyerror}.

Let $V_h \subset V$ be a conforming subspace of $V$
and throughout assume (B1) and $\psi(x,\bullet) \in C^1(\mathbb{R})$ 
for a.e.~$x \in \Omega$.
%
Given a discrete minimizer $u_h \in \arg\min E(V_h)$ of $E$ in $V_h$, let $\sigma_h \coloneqq \D W(\nabla u_h)$ denote the discrete stress. The Euler-Lagrange equations read
\begin{align}
    \int_\Omega \sigma \cdot \nabla v \d{x} &= - \int_\Omega \partial_u \psi(\bullet, u) v \d{x} \quad\text{for any } v \in V,\label{eq:ELE}\\
    \int_\Omega \sigma_h \cdot \nabla v_h \d{x} &= - \int_\Omega \partial_u \psi(\bullet, u_h) v_h \d{x} \quad\text{for any } v_h \in V_h.\label{eq:dELE-conf}
\end{align}

\begin{theorem}[a~priori of conforming methods]\label{thm:a-priori-conf}
Any $v_h \in V_h$ satisfies
    \begin{align*}
        E(u_h) - E(u) \leq \int_\Omega \big((\sigma - \sigma_h) \cdot \nabla (u - v_h) + (\partial_u \psi(\bullet,u) - \partial_u \psi(\bullet,u_h)) (u - v_h)\big) \d{x}.
\end{align*}
\end{theorem}

\begin{proof}
    The convexity of 
    $W$ and $\psi$ imply
    $0 \leq W(\nabla u) - W(\nabla u_h) - \sigma_h \cdot \nabla(u - u_h)$ 
    %a.e.~in $\Omega$ a
    and $0 \leq \psi(\bullet, u) - \psi(\bullet, u_h) - \partial_u \psi(\bullet, u_h)(u - u_h)$ a.e.~in $\Omega$. An integration provides
\begin{align}\label{ineq:pr-a-priori-conf}
        E(u_h) - E(u) &\leq
        - \int_\Omega (\sigma_h \cdot \nabla (u - u_h) + \partial_u \psi(\bullet, u_h)(u - u_h)) \d{x}.
    \end{align}
Given $v_h \in V_h$, \eqref{eq:dELE-conf} proves $(\sigma_h, v_h - u_h)_{L^2(\Omega)} = (\partial_u \psi(\bullet, u_h), u_h - v_h)_{L^2(\Omega)}$. 
This, \eqref{ineq:pr-a-priori-conf}, and 
$(\sigma, \nabla(u - v_h))_{L^2(\Omega)} 
= -(\partial_u \psi(\bullet, u), u - v_h)_{L^2(\Omega)}$ reveal
\begin{align*}
E(u_h) - E(u) &\leq - \int_\Omega (\sigma_h \cdot \nabla (u - v_h) 
+ \partial_u \psi(\bullet, u_h) (u - v_h)) \d{x}\\
&\leq \int_\Omega \big((\sigma - \sigma_h) \cdot \nabla (u - v_h) 
+ (\partial_u \psi(\bullet,u) 
- \partial_u \psi(\bullet,u_h)) (u - v_h)\big) \d{x}.
\qedhere
\end{align*}
\end{proof}

The following estimates imply \eqref{confenergyerror}.

\begin{corollary}[convergence rates of conforming methods]
Suppose (B2) and (B4). Then
    \begin{align*}
        0 \leq E(u_h) - E(u) \lesssim \min_{v_h \in V_h} \|\nabla(u - v_h)\|_p.
    \end{align*}
Suppose (B2)--(B4). Then
    \begin{align*}
        0 \leq E(u_h) - E(u) \lesssim \min_{v_h \in V_h} \|\nabla(u - v_h)\|_p^2.
    \end{align*}
\end{corollary}

\begin{proof}
    For the linear right-hand side in (B4), 
    $\partial_u \psi(\bullet,u) - \partial_u \psi(\bullet,u_h) = 0$ 
    and \Cref{thm:a-priori-conf} imply
    \begin{align}\label{ineq:a-priori-conf-lin-rhs}
        E(u_h) - E(u) \leq
        \int_\Omega (\sigma - \sigma_h) \cdot \nabla (u - v_h) \d{x} \quad\text{for any } v_h \in V_h.
    \end{align}
    The two sided growth (B2) leads to a uniform bound for $\nabla u$, $\nabla u_h$ in $L^p(\Omega)^n$ and $\sigma$, $\sigma_h$ in $L^{p'}(\Omega)^n$ \cite{GlowinskiMarrocco1975, CPlechac1997}. Hence, the first assertion follows from \eqref{ineq:a-priori-conf-lin-rhs} and the H\"older inequality.
    If (B3) holds, then the choice $\alpha \coloneqq \nabla u$ and $\beta \coloneqq \nabla u_h$ in (B3) and \eqref{eq:ELE} reveal
    \begin{align*}
        \|\sigma - \sigma_h\|^2_{L^{p'}(\Omega)} \lesssim \int_{\Omega} (W(\nabla u_h) - W(\nabla u) - \sigma \cdot (u_h - u)) \d{x} = E(u_h) - E(u).
    \end{align*}
    This, \eqref{ineq:a-priori-conf-lin-rhs}, and a Young inequality conclude the second assertion.
\end{proof}
\end{document}
